% macro conflict, first definition kept: \newcommand{\Ty}{\mathrm{Ty}}
% macro conflict, first definition kept: \newcommand{\Hom}{\underline{\mathrm{Hom}}}
% macro conflict, first definition kept: \newcommand{\Ty}{\mathrm{Ty}}
% macro conflict, first definition kept: \newcommand{\red}{\text{\upshape réd}}
\documentclass[11pt]{article}
\usepackage[no-math]{fontspec}
% Reader presentation: bilingual hyphenation and legacy-kernel compatibility.
\providecommand\IfDocumentMetadataT[1]{}
\usepackage[english,french]{babel}
\usepackage[a4paper,left=22mm,right=30mm,top=22mm,bottom=24mm,marginparwidth=19mm,marginparsep=5mm]{geometry}
\usepackage{microtype}
\setmonofont[Scale=0.95]{Latin Modern Mono}
\setlength{\emergencystretch}{2em}
\setlength{\marginparpush}{8pt}
\usepackage{amsmath}
\usepackage{amssymb}
\usepackage{amsthm}
\usepackage{mathrsfs}
\usepackage{tikz-cd}
\usepackage{enumitem}

\newcommand{\printedpage}[1]{\marginpar{\footnotesize p.~#1}}
\newcommand{\msp}[1]{\marginpar{\vspace*{10pt}\footnotesize [#1]}}
\newcommand{\uTy}{\underline{\mathrm{Ty}}}
\newcommand{\Ch}{\underline{\mathrm{Ch}}}
\newcommand{\Ty}{\underline{\mathrm{Ty}}}
\newcommand{\Nerf}{\mathrm{Nerf}}
\newcommand{\Fl}{\mathrm{Fl}}
\newcommand{\Hom}{\mathrm{Hom}}
\newcommand{\ch}{\mathrm{ch}}
\newcommand{\hot}{\mathrm{hot}}
\newcommand{\id}{\mathrm{id}}
\newcommand{\CH}{\underline{\mathrm{CH}}}
\newcommand{\red}{\mathrm{r\acute{e}d}}
\newcommand{\uHom}{\underline{\mathrm{Hom}}}
\newcommand{\ul}[1]{\underline{#1}}
\newcommand{\triv}{\mathrm{triv}}
\newcommand{\uCh}{\underline{\mathrm{Ch}}}
\newcommand{\Cat}{\mathrm{Cat}}
\newcommand{\Cathot}{\mathrm{Cat\,hot}}
\newcommand{\Hot}{\mathrm{Hot}}
\newcommand{\Ob}{\mathrm{Ob}\,}
\newcommand{\sHom}{\underline{\mathcal{H}\mathit{om}}}
\newcommand{\To}{T_{\mathrm{o}}}
\newcommand{\circled}[1]{\tikz[baseline=(c.base)]\node[draw,circle,inner sep=1pt](c){#1};}
\newcommand{\Hotz}{\mathrm{Hot}^{\wedge}_{0}}
\newcommand{\CNo}{C_{\mathbf{N}^{o}}}
\newcommand{\lab}[2][l]{\text{\footnotesize\begin{tabular}{@{}#1@{}}#2\end{tabular}}}
\begin{document}
% ---- from Der-ChVII-456_p036-040.tex
\printedpage{34}
% head: A. Grothendieck : Dérivateurs, Ch. VII, Édition des Universités de Montpellier II et Paris VII
On va décrire ces monoïdes $\mathrm{Nerf}_i(\uTy)$, et constater que ce sont tous des monoïdes \textit{libres,} tout comme $\mathrm{Nerf}_0(\uTy) = \mathrm{Ob}\,\uTy = \mathrm{Ty}$. Les applications semi-simpliciales entre ces monoïdes sont donc connus quand on les connaît sur les \textit{générateurs} de ces monoïdes libres (lesquels générateurs sont aussi les éléments indécomposables, donc définis intrinsèquement en termes de la structure multiplicative). % [sic] connus -> connues, agreement with « applications »

Pour $\mathrm{Nerf}_0(\uTy) = \mathrm{Ty}$, c'est une chose vue, les générateurs sont
\[
\Delta_+,\Delta_-,
\]
correspondant respectivement aux intervalles orientés
\[
\begin{array}{rccc}
\boldsymbol{\Delta}_+ : & 0 & \longrightarrow & 1\\[4pt]
\boldsymbol{\Delta}_- : & 0 & \longleftarrow & 1.
\end{array}
\]

\msp{page 50$'$}
Pour $\mathrm{Nerf}_1(\uTy) = \mathrm{Fl}(\uTy)$, i.e. pour les flèches, le résultat $3^o$) du paragraphe précédent (détermination de $\mathrm{Fl}(\uTy)$) s'interprète aussi en disant que $\mathrm{Fl}(\uTy)$ est le monoïde libre engendré par un système de générateurs infini
\[
\mathfrak{f}_{\varepsilon,n}\ , \qquad\quad (\varepsilon,n)\in\{+1,-1\}\times\mathbf{N},
\]
indexé par l'ensemble produit $\{+1,-1\}\times\mathbf{N}$. Ici $\mathfrak{f}_{\varepsilon,n}$ désigne l'unique flèche dans $\uTy$
\[
\mathfrak{f}_{\varepsilon,l} : \Delta_\varepsilon \longrightarrow \Delta^l_\varepsilon
\]
(où bien sûr l'exposant $l$ désigne une exponentiation pour la loi de composition dans $\uTy$
\[
\Delta^l_\varepsilon = \underbrace{\Delta_\varepsilon\circ\Delta_\varepsilon\circ\ \cdots\ \circ\Delta_\varepsilon}_{l\text{ copies}},
\]
et non la puissance cartésienne dans Cat, où dans $\uTy$ (où elle n'existe que si $l\in\{0,1\}$)). % [sic] où dans -> ou dans, conjunction intended
Dans la description de $\mathrm{Fl}(\uTy)$ dans le paragraphe précédent, on a écrit le `terme général' du monoïde libre en groupant ensemble les groupes de foncteurs successifs de la forme $f_{+,0}$, $f_{-,0}$, i.e. corréspondant aux deux flèches de dégénérescence indécomposables (pour $*$) % [sic] corréspondant -> correspondant, accent
\[
\left\{
\begin{array}{lccc}
f_{+,0} : & \Delta_+ & \longrightarrow & e\\[4pt]
f_{-,0} : & \Delta_- & \longrightarrow & e
\end{array}
\right.
\]
(\footnote[15]{\textbf{NB} Le `mot' vide de $\mathrm{Fl}(\uTy)$ correspond à la flèche $\mathrm{id}_1$, où $1$ est l'élément unité de Ty, réalisant l'`intervalle' ponctuel.}). Le sous-monoïde libre de $\mathrm{Fl}(\uTy)$ engendré par ces deux éléments est le monoïde des flèches de dégénérescence (pure) quelconques
% foot: Version 22 décembre 2001
\printedpage{35}
% head: A. Grothendieck : Dérivateurs, Ch. VII, Édition des Universités de Montpellier II et Paris VII
\msp{page 51}
dans $\uTy$ :
\[
f_\tau : I_\tau \longrightarrow e \qquad\quad (\tau\in\mathrm{Ty}),
\]
et on a
\[
f_\tau * f_{\tau'} = f_{\tau\tau'},
\]
i.e. l'application
\[
\tau \longmapsto f_\tau :\quad \underbrace{\mathrm{Ob}\,(\uTy)}_{=\,\mathrm{Ty}\,=\,\mathrm{Nerf}_0(\uTy)} \quad\longrightarrow\quad \underbrace{\mathrm{Fl}(\uTy)}_{\mathrm{Nerf}_1(\uTy)}
\]
est un homomorphisme de monoïdes, transformant $\Delta_\varepsilon$ ($\varepsilon\in\{+1,-1\}$) dans $\mathfrak{f}_{\varepsilon,0}$. Ce n'est pas l'application de dégénérenscence de $\mathrm{Nerf}_0$ dans $\mathrm{Nerf}_1$, i.e. l'application $\tau\longmapsto\mathrm{id}_\tau$. Celle-ci est donnée par % [sic] dégénérenscence -> dégénérescence, spelling
\[
\begin{array}{ccc}
\mathrm{Ob}\,(\uTy) & \longrightarrow & \mathrm{Fl}(\uTy)\\[4pt]
\tau & \longmapsto & \mathrm{id}_\tau
\end{array}
\ :\
\left|\
\begin{array}{rcccl}
 & \Delta_\varepsilon & \longmapsto & \mathfrak{f}_{\varepsilon,1} & \qquad(\varepsilon\in\{+1,-1\})\\[4pt]
\text{i.e.} & \Delta_+ & \longmapsto & \mathfrak{f}_{+,1} & \\[4pt]
 & \Delta_- & \longmapsto & \mathfrak{f}_{-,1}\ . &
\end{array}
\right.
\]
Quant aux homomorphismes source et but
\[
\begin{tikzcd}[column sep=large]
\mathrm{Fl}(\uTy) \arrow[r, shift left, "\sigma\text{ (source)}"] \arrow[r, shift right, "\beta\text{ (but)}"'] & \mathrm{Ob}\,(\uTy),
\end{tikzcd}
\]
ils sont évidents sur les générateurs
\[
\left\{
\begin{array}{lcl}
\sigma(\mathfrak{f}_{\varepsilon,l}) & = & \Delta_\varepsilon\\[4pt]
\beta(\mathfrak{f}_{\varepsilon,l}) & = & \Delta^l_\varepsilon.
\end{array}
\right.
\]

Je vais décrire maintenant $\mathrm{Nerf}_2(\uTy)$, i.e. les systèmes
\[
\tau \xrightarrow{\ f\ } \tau' \xrightarrow{\ g\ } \tau''.
\]

\msp{page 52}
Se donner $f\in\mathrm{Fl}(\uTy)$ revient à se donner une suite ordonnée de générateurs $f_{\varepsilon,l}$, donc une suite
\[
\begin{aligned}
&\mathfrak{f}_{\varepsilon_1,l_1},\ \mathfrak{f}_{\varepsilon_2,l_2},\ \ \dots\ \ ,\mathfrak{f}_{\varepsilon_n,l_n}\\
&((\varepsilon_i,l_i))_{1\leq i\leq n}\in[\{\pm1\}\times\mathbf{N}]^n.
\end{aligned}
\]
On a alors
\[
\begin{aligned}
x &= \Delta_{\varepsilon_1}\Delta_{\varepsilon_2}\cdots\Delta_{\varepsilon_n} &\qquad\qquad& \mathrm{long}\,x = n\\
y &= \Delta^{l_1}_{\varepsilon_1}\Delta^{l_2}_{\varepsilon_2}\cdots\Delta^{l_n}_{\varepsilon_n} && \mathrm{long}\,y = \textstyle\sum_{i\in[1,n]} l_i.
\end{aligned}
\]
% foot: Version 22 décembre 2001

\printedpage{36}
% head: A. Grothendieck : Dérivateurs, Ch. VII, Édition des Universités de Montpellier II et Paris VII
Se donner $q : y\longrightarrow z$, i.e. un objet de $y\backslash\uTy$, revient à se donner une application de l'ensemble des $m=\sum l_i$ flèches de transition de $y$ dans $\mathbf{N}$. En fait, on a % [sic] q -> g, the map is called g below (« la donnée de g »)
\[
\mathrm{Tr}(y) = \mathrm{Tr}(\Delta^{l_1}_{\varepsilon_1})\amalg\mathrm{Tr}(\Delta^{l_2}_{\varepsilon_2})\amalg\ \cdots\ \amalg\mathrm{Tr}(\Delta^{l_n}_{\varepsilon_n})
\]
(`ensemble somme', i.e. réunion disjointe), donc
\[
\mathrm{Hom}_{\mathrm{Ens}}(\mathrm{Tr}(y)) \simeq \prod_{i\in[1,n]}\mathrm{Hom}_{\mathrm{Ens}}(\underbrace{\mathrm{Tr}(\Delta^{l_i}_{\varepsilon_i})}_{[1,l_i]},\mathbf{N}), % [sic] Hom_Ens(Tr(y)) -> Hom_Ens(Tr(y),N), second argument missing
\]
donc la donnée de $g$ équivaut à la donnée d'un système de $n$ applications
\[
\left\{
\begin{array}{l}
\pi_i : [1,l_i]\longrightarrow\mathbf{N} \qquad\quad (\text{pondération de }\Delta^{l_i}_{\varepsilon_i})\\[6pt]
\text{pour } i\in[1,n].
\end{array}
\right.
\]
Mais pour $i\in[1,n]$, considérons l'élément de $\mathrm{Nerf}_2$
\[
\Delta_{\varepsilon_1}\xrightarrow{\mathfrak{f}_{\varepsilon_i,l_i}}\Delta^{l_i}_{\varepsilon_i}\xrightarrow{g_{\varepsilon,\pi_i}}\underbrace{\Delta^{\pi_i(1)}_{\varepsilon_i}\Delta^{\pi_i(2)}_{\varepsilon_i}\cdots\Delta^{\pi_i(l_i)}_{\varepsilon_i}}_{\Delta_{\varepsilon_i}^{\Sigma_j\pi_i(j)\overset{\text{déf}}{=}\mathrm{long}(\pi_i)}}, % [sic] \Delta_{\varepsilon_1} -> \Delta_{\varepsilon_i}, and g_{\varepsilon,\pi_i} -> g_{\varepsilon_i,\pi_i}, index i intended
\tag*{(1)}
\]

\msp{page 53}
où $g_{\varepsilon_i,\pi_i}$ est la flèche définie en termes de $\pi_i$ au paragraphe précédent. Mais ce n'est autre aussi que
\[
g_{\varepsilon_i,\pi_i} = \mathfrak{f}_{\varepsilon_i,\pi_i(1)}*\mathfrak{f}_{\varepsilon_i,\pi_i(2)}*\cdots*\mathfrak{f}_{\varepsilon_i,\pi_i(l_i)}.
\tag*{(2)}
\]
L'élément de $\mathrm{Nerf}_2$ envisagé ne dépend que de
\[
(\varepsilon_i,l_i,\pi_i), \qquad\quad \text{où } (\varepsilon_i,l_i)\in\{\pm1\}\times\mathbf{N},\ \pi_i : [1,l_i]\longrightarrow\mathbf{N}.
\]
(\textbf{NB} Si $l_i=0$, $[1,l_i]=[1,0]=\emptyset$ par convention, et $\pi_i$ est l'unique application $\emptyset\longrightarrow\mathbf{N}$, donc la donnée $(\varepsilon_i,l_i,\pi_i)$ revient alors à $(\varepsilon_i,l_i)$.)

Si
\[
(\varepsilon,l,\pi)\ \text{ comme dessus,} \qquad\quad (\varepsilon,l)\in\{\pm1\}\times\mathbf{N},\ \pi : [1,l]\longrightarrow\mathbf{N},
\]
désignons par
\[
\mathfrak{f}_{\varepsilon,l,\pi}
\]
l'élément de $\mathrm{Nerf}_2$ donné par
\[
\Delta_\varepsilon\xrightarrow{\mathfrak{f}_{\varepsilon,l}}\Delta^l_\varepsilon\ \xrightarrow{\quad\mathfrak{f}_{\varepsilon,\pi(1)}\ \ *\ \cdots\ *\ \ \mathfrak{f}_{\varepsilon,\pi(l)}\quad}\Delta^{\pi(1)}_\varepsilon\Delta^{\pi(2)}_\varepsilon\cdots\Delta^{\pi(l)}_\varepsilon.
\]
On peut formuler la construction précédente en disant que tout élément de $\mathrm{Nerf}_2$ peut se décomposer de façon unique en un produit
\[
(f,g) = \mathfrak{f}_{\varepsilon_1,l_1,\pi_1}*\mathfrak{f}_{\varepsilon_2,l_2,\pi_2}*\cdots*\mathfrak{f}_{\varepsilon_n,l_n\pi_n}, % [sic] l_n\pi_n -> l_n,\pi_n, comma missing in index
\]
ce qui prouve que $\mathrm{Nerf}_2(\uTy)$ est le monoïde libre sur les générateurs précédents $\mathfrak{f}_{\varepsilon,l,\pi}$. (\textbf{NB} à la suite vide correspond
\[
1\xrightarrow{\mathrm{id}_1}1\xrightarrow{\mathrm{id}_1}1 \qquad\qquad .)
\]
% foot: Version 22 décembre 2001

\printedpage{37}
% head: A. Grothendieck : Dérivateurs, Ch. VII, Édition des Universités de Montpellier II et Paris VII
\msp{page 54}
Si l'objet de $\mathrm{Nerf}_2$ est
\[
x\xrightarrow{\ f\ }y\xrightarrow{\ g\ }z,
\]
la construction précédente nous dit, en termes de la donnée combinatoire, i.e. de la suite de données
\[
(\varepsilon_1,l_1,\pi_1),\ \dots\ ,\ (\varepsilon_n,l_n,\pi_n),
\]
quels sont $x$, $y$, $z$ : chacun d'eux est le produit dans $\mathrm{Nerf}_0=\mathrm{Ty}$ des $n$ facteurs $x_i$, $y_i$ ou $z_i$ respectivement, pour $i\in[1,n]$, donnés dans la définition des $\mathfrak{f}_{\varepsilon_i,l_i,\pi_i}$, i.e. de
\[
\Delta_{\varepsilon_i},\ \Delta^{l_i}_{\varepsilon_i},\ \text{ et }\Delta^{\pi_i(1)}_{\varepsilon_i}\Delta^{\pi_i(2)}_{\varepsilon_i}\cdots\Delta^{\pi_i(l_i)}_{\varepsilon_i}=\Delta^{\mathrm{long}(\pi_i)}_{\varepsilon_i}
\]
(où $\mathrm{long}(\pi_i)=\sum_{j\in[1,l_i]}\pi_i(j)$). Donc
\[
\left\{
\begin{array}{lcl}
x & = & \prod^{\mathrm{Ty}}_{i\in[1,n]}\Delta_{\varepsilon_i}\\[4pt]
y & = & \prod^{\mathrm{Ty}}_{i\in[1,n]}\Delta^{l_i}_{\varepsilon_i}\\[4pt]
z & = & \prod^{\mathrm{Ty}}_{i\in[1,n]}\Delta^{\mathrm{long}(\pi_i)}_{\varepsilon_i}
\end{array}
\right.
\]
(\footnote[16]{attention, les produits doivent être pris en respectant l'ordre des facteurs donné par la notation indicielle en $i$.}), et on trouve de même $f$, $g$ comme des produits (au sens $*$) de $n$ facteurs $f_i$ resp. $g_i$ ($i\in[1,n]$) explicités dans les formules (1), (2) (p. 52, 53) :
\[
\begin{array}{lclccl}
f_i & = & \mathfrak{f}_{\varepsilon_i,l_i} : & \Delta_{\varepsilon_i} & \longrightarrow & \Delta^{l_i}_{\varepsilon_i}\\[4pt]
g_i & = & g_{\varepsilon_i,l_i} : & \Delta^{l_i}_{\varepsilon_i} & \longrightarrow & \Delta^{\mathrm{long}(\pi_i)}_{\varepsilon_i}, % [sic] g_{\varepsilon_i,l_i} -> g_{\varepsilon_i,\pi_i}, notation of (1), (2)
\end{array}
\]
$g_i$ lui-même étant représenté comme un $*$-produit de $l_i$ facteurs indécomposables, de sorte que $f$ et $g$ apparaissent sous forme
\msp{page 55}
de $*$-produits de $n$, resp. de $\sum_{i\in[1,n]}l_i$ facteurs de la forme $\mathfrak{f}_{\varepsilon,\nu}$. Pour terminer de tenir en mains l'objet semi-simplicial tronqué $(\mathrm{Nerf}_i(\uTy))_{0\leq i\leq 2}$, il reste à décrire les deux applications de dégénérescence $\mathrm{Nerf}_1\longrightarrow\mathrm{Nerf}_2$ (cas où $f$ ou $g$ est une flèche identique) et surtout l'application de \textit{composition} des flèches
\[
\mathrm{Nerf}_2(\uTy)\xrightarrow{\ \gamma\ }\mathrm{Nerf}_1(\uTy)=\mathrm{Fl}(\uTy).
\]
Il suffit de donner $\gamma$ sur les générateurs, et on a tautologiquement, par définition des $\mathfrak{f}_{\varepsilon,l,\pi}$
\[
\gamma(\mathfrak{f}_{\varepsilon,l,\pi})=\mathfrak{f}_{\varepsilon,\mathrm{long}(\pi)}.
\]
(\textbf{NB} Il n'existe au plus qu'un seul morphisme de $\Delta_\varepsilon$ (de $\Delta_+$ ou de $\Delta_-$) dans un intervalle $I_\tau$. Il existe si et seulement si $\tau$ est de la forme $\Delta^l_\varepsilon$, $l\in\mathbf{N}$.)
% foot: Version 22 décembre 2001

\printedpage{38}
% head: A. Grothendieck : Dérivateurs, Ch. VII, Édition des Universités de Montpellier II et Paris VII
Le cas $g=\mathrm{id}$ correspond au cas où
\[
\pi_i(j)=1 \qquad\quad \forall j\in[1,l_i],
\]
les $\mathfrak{f}_{\varepsilon_i,l_i,\pi_i}$ sont alors détérminés en fonction des $\mathfrak{f}_{\varepsilon_i,l_i}$, et l'application de dégénérescence est % [sic] détérminés -> déterminés, accent
\[
\mathfrak{f}_{\varepsilon,l}\longmapsto\mathfrak{f}_{\varepsilon,l,1_l} \qquad\quad (\text{où } 1_l \text{ est la fonction constante de valeur 1 sur } [1,l]).
\]
Le $f=\mathrm{id}$ est celui où les $l_i$ sont égaux à 1. Alors chaque $\pi_i : [1,1]\longrightarrow\mathbf{N}$ peut être vu comme un élément de $\mathbf{N}$, la donnée de $(\varepsilon_i,l_i,\pi_i)$ revient à celle des $(\varepsilon_i,n_i)\in\{\pm1\}\times\mathbf{N}$, i.e. la donnée % [sic] Le f = id -> Le cas f = id, word « cas » missing
\msp{page 56}
d'une suite de $n$ générateurs parmi les $\mathfrak{f}_{\varepsilon,l}$. L'application de dégénérescence $\mathrm{Nerf}_1\longrightarrow\mathrm{Nerf}_2$ est ici donnée sur les générateurs par
\[
\mathfrak{f}_{\varepsilon,l}\longmapsto\mathfrak{f}_{\varepsilon,1,l}.
\]
Heureusement, on n'a plus à prouver l'associativité de la composition.

\textbf{NB} La simplicité de ces descriptions renforce mon sentiment que la catégorie $\uTy$ définie à présent, avec les \textit{morphismes d'intervalles orientés croissants} $I_\tau\longrightarrow I_{\tau'}$ comme flèches de $\uTy$, est bien `la bonne'. On va déjà voir si elle donne une description sympa des catégories de fractions $X\Sigma^{-1}$, sans tortillements comme précédemment.

Mais on a promis de déterminer $\mathrm{Nerf}_r(\uTy)$, pour tout $r\in\mathbf{N}$. Le générateur général de $\mathrm{Nerf}_r$ est de la forme (pour $\varepsilon\in\{\pm1\}$ fixé)
\[
\underbrace{\Delta_\varepsilon}_{=\,\Delta^{l_0}_\varepsilon\text{avec } l_0=1}\xrightarrow{\ f_1\ }\Delta^{l_1}_\varepsilon\xrightarrow{\ f_2\ }\Delta^{l_2}_\varepsilon\longrightarrow\ \cdots\ \xrightarrow{\ f_r\ }\Delta^{l_r}_\varepsilon,
\]
où les $f_i : \Delta^{l_{i-1}}_\varepsilon\longrightarrow\Delta^{l_i}_\varepsilon$, $i\in[1,r]$, sont donnés par des applications
\[
\pi_{i-1} : [1,l_{i-1}]\longrightarrow\mathbf{N} \qquad\quad \text{telles que }\ \underbrace{\mathrm{long}(\pi_{i-1})}_{\overset{\text{déf}}{=}\ \Sigma_{j\in[1,l_i]}\pi_i(j)}\ = l_i. % [sic] \Sigma_{j\in[1,l_i]}\pi_i(j) -> \Sigma_{j\in[1,l_{i-1}]}\pi_{i-1}(j), indices shifted with respect to long(\pi_{i-1})
\]
(NB $\pi_0 : [1,1]=\{1\}\longrightarrow\mathbf{N}$ est identifié à $\mathrm{long}(\pi_0)=\pi_0(1)=l_1$.)

\msp{page 57}
Pour ne pas faire dépendre la définition du domaine des $\pi_i$ ($i\in[0,r-1]$) des entiers $l_i$, on peut aussi considérer les $\pi_i$ comme des applications
\[
\mathbf{N}^*\xrightarrow{\ \pi\ }\mathbf{N} \qquad\quad (\mathbf{N}^*=\{i\in\mathbf{N}\mid i\geq 1\})
\]
à supports finis (i.e. tels que $\pi(i)=0$ pour $i$ grand), et exiger de façon précise
\[
\begin{aligned}
&\mathrm{supp}\,\pi_0\ \subseteq\ \{1\}\\
&\mathrm{supp}\,\pi_1\ \subseteq\ [1,\mathrm{long}(\pi_0)]\\
&\qquad\ \cdots\\
&\boxed{\mathrm{supp}\,\pi_i\ \subseteq\ [1,\mathrm{long}(\pi_{i-1})]}, \qquad\quad \text{pour }\ \boxed{i\in[0,r-1]}\quad .
\end{aligned}
\tag*{$(**)$}
\]
% foot: Version 22 décembre 2001

% ---- from Der-ChVII-456_p041-045.tex
\printedpage{39}
% head: A. Grothendieck : Dérivateurs, Ch. VII, Édition des Universités de Montpellier II et Paris VII
D'autre part, posant $l_i = \mathrm{long}(\pi_{i-1})$ (pour $i \in [1,r]$),
\[
f_i : \Delta_\varepsilon^{l_{i-1}} \longrightarrow \Delta_\varepsilon^{l_i\ (=\mathrm{long}(\pi_{i-1}))}
\]
est défini comme
\[
f_i = \mathfrak{f}_{\varepsilon,\pi_{i-1}(1)} * \mathfrak{f}_{\varepsilon,\pi_{i-1}(2)} * \cdots * \mathfrak{f}_{\varepsilon,\pi_{i-1}(l_{i-1})},
\]
de sorte que l'élément $(*)$ de $\Nerf_r(\Ty)$ est déterminé de façon unique par la donnée des $\varepsilon \in \{\pm 1\}$ et de
\[
\pi_\bullet = (\pi_i)_{i\in[0,r-1]} \in (\mathrm{Hom}_{\mathrm{Ens}}(\mathbf{N}^*,\mathbf{N}))^r
\]
satisfaisant les conditions $(**)$. Soit $F_{\varepsilon,\pi_\bullet}$ cet élément de $\Nerf_r(\Ty)$, et soit $N_r$ l'ensemble des $\pi_\bullet$ qu'on vient d'expliciter. Alors
\[
(F_{\varepsilon,\pi_\bullet})_{(\varepsilon,\pi_\bullet)\in\{\pm1\}\times N_r}
\]
est un système de générateurs libre du monoïde $\Nerf_r(\Ty)$.

Si $l_i = 0$ (i.e. \fbox{$\pi_{i-1} = 0$}), on doit avoir $\mathrm{supp}(\pi_{i+1}) = \emptyset$, i.e. $\pi_{i+1} = 0$, d'où $l_{i+1} = 0$, et de même pour les $l_j$ et $\pi_j$ suivants. Géométriquement l'hypothèse s'exprime par la condition que le but
\msp{page 58}
de $f_i$ est l'élément 1 de Ty (correspondant à l'intervalle ponctuel), et comme $\mathrm{Hom}(1,\tau) = \emptyset$ pour $\tau \neq 1$, cela signifie que les flèches qui suivent $f_i$ doivent être toutes égales à la flèche $\mathrm{id}_1 : 1 \longrightarrow 1$.

\textbf{Remarque.} Dans la description combinatoire de $\Ty$ et de $\Nerf_\bullet(\Ty)$, l'ensemble des `signatures' $\{\pm1\}$ joue un rôle assez bidon. Il peut être remplacé par n'importe quel ensemble $\Sigma$ (dit `ensemble des signatures'), et on construit alors une catégorie monoïdale (ou un monoïde catégorique) ayant comme ensemble d'objets le monoïde libre engendré par $\Sigma$, interprétés comme des intervalles d'entiers $[0,n]$, munis de `signatures' dans $\Sigma$ pour chaque intervalle de transition, ou ce qui revient au même, des `graphes' orientés $\Sigma$-pondérés de type combinatoire un intervalle $\boldsymbol{\Delta}_i$, ce qui revient à la donnée d'une suite (ordonnée) de $n$ éléments $\varepsilon_1,\ \varepsilon_2,\ \ldots\ ,\ \varepsilon_n \in \Sigma$, i.e. d'un `mot' $\varepsilon_1\varepsilon_2\ldots\varepsilon_n$, qu'on préfère écrire $\Delta_{\varepsilon_1}\Delta_{\varepsilon_2}\cdots\Delta_{\varepsilon_n}$. On peut définir de façon évidente les homomorphismes \textit{croissants} de tels graphes pondérés orientés (sous-entendu : avec l'élément source et l'élément but dans la structure, à respecter par les morphismes), sans avoir besoin d'une interprétation catégorique. La fonctorialité de l'opération de composition des intervalles pondérés, pour cette notion de morphisme, est évidente. Le
\msp{page 59}
choix $\Sigma = \{\pm1\}$ est à ma connaissance le seul utile, à cause de son lien avec la notion de chemins et le calcul de $\pi_0$ d'un objet de $X$ \texttt{[plutôt d'un objet} $X$ \texttt{de} Cat\texttt{]}, etc.
% foot: Version 22 décembre 2001

\printedpage{40}
% head: A. Grothendieck : Dérivateurs, Ch. VII, Édition des Universités de Montpellier II et Paris VII
\section*{7\quad Explicitation des catégories de fractions $X\Sigma^{-1}$ en termes des $\Ch(X;a,b)$}

Si $X$ est une catégorie, on désigne par $\Ch(X)$ la catégorie des chemins dans $X$ (cf. §4), et pour $a,b \in \mathrm{Ob}\,X$, par
\[
\Ch(X;a,b) \qquad \text{ou simplement } \Ch(a,b)
\]
celle des chemins de $a$ vers $b$, i.e. de source $a$, but $b$. L'intuition géométrique (qu'on explicitera), c'est qu'un tel chemin $c$ définit un chemin `direct' de $a$ vers $b$ (i.e. un chemin de type $\tau'$ `direct', i.e. toutes les flèches de transition sont directes, i.e. $\tau'$ de la forme $\Delta_+^n$), de même longueur que le chemin initial $c$, obtenu en inversant toutes les flèches rétrogrades dans le chemin $c$. Cette opération est effectivement possible dans $X$ lui-même si et seulement si lesdites flèches sont inversibles. On obtient alors une suite composable de flèches, dont le composé total (qu'on peut appeler `l'effet utile' du chemin)
\msp{page 60}
est une flèche de $a$ vers $b$. Dans le cas général, les chemins de $a$ vers $b$ peuvent être vus comme une sorte de `flèches généralisées' allant de $a$ vers $b$, dont la composition serait la composition évidente, et les identités les `chemins vides' relatifs aux $a \in \mathrm{Ob}\,X$. On trouve d'ailleurs bien ainsi une catégorie, ayant mêmes objets que $X$, mais où les
\[
\mathrm{Hom}_{\mathrm{Ch}}(a,b) = \mathrm{Ch}(a,b) = \mathrm{Ob}\,\Ch(a,b)
\]
sont prohibitivement gros. Cela s'exprime notamment par le fait que les seules flèches inversibles dans $X_{\mathrm{Ch}}$ sont les identités ! On n'a pas de foncteur naturel entre $X_{\mathrm{Ch}}$ et $X$ qui soit l'identité sur les objets. Car à un chemin $a \xrightarrow{c} b$ on ne peut associer de $X$-morphismes $a \longrightarrow b$ (sauf si $c$ \texttt{[est]} de type $\Delta_+$), et à un $X$-morphisme $a \xrightarrow{u} b$ on peut associer un chemin $c(u)$ de type $\Delta_+$, mais si $a \xrightarrow{u} b \xrightarrow{v} c$, $c(vu)$ est encore de type $\Delta_+$ alors que $c(v)c(u)$ est de type $\Delta_+^2$, ils ne sont pas égaux. Par contre, il y a une flèche naturelle entre ces deux éléments de $\Ch(a,c)$, flèche au dessus de l'unique homomorphisme
\[
\Delta_+ \xrightarrow{\mathfrak{f}_{+,2}} \Delta_+^2.
\]
\msp{page 61}
On a alors
\[
\mathfrak{f}_{+,2}^*(\underbrace{a \xrightarrow{u} b \xrightarrow{v} c}_{c(v)c(u)}) = (\underbrace{a \xrightarrow{vu} c}_{c(vu)}),
\]
d'où une flèche canonique
\[
c(vu) \longrightarrow c(v)c(u) \qquad \text{dans } \Ch(a,c).
\]
Cela montre que $c(vu)$ et $c(v)c(u)$ sont tout au moins dans une même composante connexe de $\Ch(a,c)$, ou (comme on dit encore) qu'ils sont `homotopes'. Cela suggère ici de modifier la définition de $\mathrm{Hom}_{\mathrm{Ch}}(a,b)$ en
\[
\pi(a,b) = \pi_0(\Ch(a,b)).
\]
% foot: Version 22 décembre 2001

\printedpage{41}
% head: A. Grothendieck : Dérivateurs, Ch. VII, Édition des Universités de Montpellier II et Paris VII
Le foncteur de composition des chemins
\[
\begin{array}{ccccc}
\Ch(a,b) & \times & \Ch(b,c) & \longrightarrow & \Ch(a,c)\\
(c & , & c') & \longmapsto & c' \circ c
\end{array}
\]
\texttt{[il y a trop de} $c$\texttt{]} (\footnote[17]{Ce foncteur est déduit des foncteurs
\[
\Ch_\tau(a,b) \times \Ch_{\tau'}(b,c) \longrightarrow \Ch_{\tau\tau'}(a,c)
\]
évidents, cf. page 77.}) induit une application sur les $\pi_0$, ce qui donne une loi de composition
\[
\pi(a,b) \times \pi(b,c) \longrightarrow \pi(a,c)
\]
sur les flèches. Il est immédiat qu'elle est associative (la composition des chemins l'étant), et unitaire. Elle admet comme flèche identique en $x$ soit le `lacet vide en $x$', lequel est (directement) homotope dans $\Ch(x,x)$ au chemin `dégénéré' $x \xrightarrow{\mathrm{id}} x$
\msp{page 62}
de type $\Delta_+$, ou aussi au chemin dégénéré de n'importe quel type $\tau$, disons
\[
x \xleftarrow{\mathrm{id}} x \xrightarrow{\mathrm{id}} x \xrightarrow{\mathrm{id}} x \xleftarrow{\mathrm{id}} x \xleftarrow{\mathrm{id}} x,
\]
un tel chemin n'étant autre que $p_\tau^*(x)$, où
\[
p_\tau : \tau \longrightarrow 1
\]
est le morphisme de dégénérescence (l'unique morphisme dans $\Ty$ vers 1, qui est objet final), d'où
\[
p_\tau^* : \Ch_1(x,x) = \{x\} \longrightarrow \Ch_\tau(x,x) \subseteq \Ch(x,x).
\]
La catégorie $\Pi(X)$ ainsi obtenu doit s'appeler le \textit{groupoïde fondamental} de $X$. En effet, % [sic] obtenu -> obtenue, agreement with « La catégorie »

\begin{enumerate}[label=\alph*)]
\item toutes les flèches de $\Pi(X)$ sont inversibles, i.e. $\Pi(X)$ est bien un groupoïde, et
\item le foncteur canonique
\[
X \longrightarrow \Pi(X)
\]
qui est l'identité sur les objets, et qui est défini sur les flèches par $u \longmapsto c(u)$ comme ci-dessus, est \textit{universel} (dans Cat) pour les foncteurs (i.e. les Cat-morphismes) de $X$ vers des \textit{groupoïdes.} (Plus généralement encore, il est universel pour les foncteurs $X \xrightarrow{F} Y$ tels que $F$ transforme toute flèche de $X$ en \texttt{[une]} flèche inversible.)
\end{enumerate}
% foot: Version 22 décembre 2001

\printedpage{42}
% head: A. Grothendieck : Dérivateurs, Ch. VII, Édition des Universités de Montpellier II et Paris VII
\msp{page 63}
Pour prouver que toute flèche dans $\Pi(X)$ est inversible, comme tout chemin non vide est composé de chemins `élémentaires' de type $\Delta_+$ ou $\Delta_-$, il suffit de prouver que ces chemins sont inversibles. En fait, ils se correspondent biunivoquement,
\[
\Ch_{\Delta_+}(a,b) \simeq \Ch_{\Delta_-}(b,a),
\]
en associant au chemin de $a$ vers $b$
\[
c(u) : a \xrightarrow{u} b \qquad \text{de type } \Delta_+
\]
le chemin
\[
c(u)^o : b \xleftarrow{u} a \qquad \text{de type } \Delta_-
\]
de $b$ vers $a$. Je dis que
\[
\begin{array}{rcl}
c(u)^o c(u) & \sim & 1_a\\
c(u) c(u)^o & \sim & 1_b
\end{array}
\]
\[
\left.
\begin{tikzcd}
 & a \arrow[r, "u"] & b & a \arrow[l, "u"'] \\
 {} & a \arrow[u, "\mathrm{id}"'] \arrow[r, "\mathrm{id}"] & a \arrow[u, "u"'] & a \arrow[u, "\mathrm{id}"'] \arrow[l, "\mathrm{id}"'] \\
 {} \arrow[u] & & a &
\end{tikzcd}
\right\}\
\begin{array}{l}
\text{on a un homomorphisme dans } \Ch_{\Delta_+\Delta_-}(a,a)\\
\text{de } 1_a \text{ dans } c(u)^o c(u).
\end{array}
\]
\[
\left.
\begin{tikzcd}
 & b & a \arrow[l, "u"'] \arrow[r, "u"] & b \\
 {} & a \arrow[u, "u"'] & a \arrow[u, "\mathrm{id}"'] \arrow[l, "\mathrm{id}"'] \arrow[r, "\mathrm{id}"] & a \arrow[u, "u"'] \\
 {} \arrow[u] & & a &
\end{tikzcd}
\right\}\
\begin{array}{l}
\text{itou dans } \Ch_{\Delta_-\Delta_+}(b,b),\\
\text{de } 1_b \longrightarrow c(u)c(u)^0 % [sic] c(u)^0 -> c(u)^o, digit 0 printed instead of the letter o used elsewhere
\end{array}
\]
\texttt{[plutôt} $1_b \longleftarrow c(u)c(u)^0$\texttt{, viz.} % [sic] c(u)^0 -> c(u)^o, digit 0 printed instead of the letter o used elsewhere
\[
\begin{tikzcd}
b \arrow[d, "\mathrm{id}"] & a \arrow[l, "u"'] \arrow[r, "u"] \arrow[d, "u"] & b \arrow[d, "\mathrm{id}"] \\
b & b \arrow[l, "\mathrm{id}"'] \arrow[r, "\mathrm{id}"] & b \ \texttt{].}
\end{tikzcd}
\]
% foot: Version 22 décembre 2001

\printedpage{43}
% head: A. Grothendieck : Dérivateurs, Ch. VII, Édition des Universités de Montpellier II et Paris VII
\msp{page 64}
Notons que, plus généralement, pour tout chemin
\[
c : a \longrightarrow b, \qquad \text{de type } \tau,
\]
on définit un \textit{chemin opposé}
\[
c^o : b \longrightarrow a \qquad \text{de type } \tau^o,
\]
où
\[
\tau \longmapsto \tau^o : \mathrm{Ty} \longrightarrow \mathrm{Ty}
\]
est l'antiïnvolution du monoide libre Ty qui échange les deux générateurs $\Delta_-$, $\Delta_+$, et qui est donc donné sur l'élément général par
\[
(\Delta_{\varepsilon_1}\Delta_{\varepsilon_2}\cdots\Delta_{\varepsilon_n})^o = \Delta_{-\varepsilon_n}\Delta_{-\varepsilon_{n-1}}\cdots\Delta_{-\varepsilon_1}.
\]
Par récurrence sur la longueur de $c$ (\footnote[18]{inutile, résulte directement de la formule $(*)$ \texttt{[de la]} page suivante, en décomposant $c$ en composés de chemins `élémentaires' de type $\Delta_+$ ou $\Delta_-$.}), en utilisant le résultat précédent, que $c(u)$, $c(u)^o$ sont inverses l'un de l'autre dans $\Pi(X)$, on trouve
\[
\left.
\begin{array}{rcl}
c^o c & = & \mathrm{id}_a\\
c c^o & = & \mathrm{id}_b
\end{array}
\right\}\text{dans } \Pi(X).
\]

\textbf{Exemple.} Si
\[
\begin{array}{rcl}
\tau & = & \Delta_+\Delta_-\Delta_+^2\Delta_-\\
\tau^o & = & \Delta_+\Delta_-^2\Delta_+\Delta_-,
\end{array}
\]
et si on a un chemin
\[
c : a_0 \xrightarrow{u_1} a_1 \xleftarrow{u_2} a_2 \xrightarrow{u_3} a_3 \xrightarrow{u_4} a_4 \xleftarrow{u_5} a_5
\]
de type $\tau$, on en déduit le chemin
\[
c^o : a_5 \xrightarrow{u_5} a_4 \xleftarrow{u_4} a_3 \xleftarrow{u_3} a_2 \xrightarrow{u_2} a_1 \xleftarrow{u_1} a_0
\]
de type $\tau^o$. L'opération $c \longmapsto c^o$ satisfait en
\msp{page 65}
effet \texttt{[pour]}
\[
\begin{array}{ccccc}
a & \xrightarrow{\gamma} & b & \xrightarrow{\gamma'} & c\\
c & \xrightarrow{\gamma'^o} & b & \xrightarrow{\gamma^o} & a
\end{array}
\]
la relation
\[
(\gamma'\gamma)^o = \gamma^o\gamma'^o, \tag*{$(*)$}
\]
% foot: Version 22 décembre 2001

% ---- from Der-ChVII-456_p046-050.tex
\printedpage{44}
% head: A. Grothendieck : Dérivateurs, Ch. VII, Édition des Universités de Montpellier II et Paris VII
% foot: Version 22 décembre 2001

ce qui est en accord avec l'intuition que l'opération $\gamma \longmapsto \gamma^o$ joue le rôle d'un passage à l'inverse (`à homotopie près'), et on a la formule
\[
(\gamma'\gamma)^{-1} = \gamma^{-1}\gamma'^{-1}
\]
pour deux flèches inversibles $\gamma$, $\gamma'$ composables (dans l'ordre $\gamma'$, $\gamma$).

Pour la propriété universelle de
\[
X \longrightarrow \Pi(X)
\]
(qui permettra d'écrire
\[
\Pi(X) \simeq X\Sigma^{-1}, \qquad \text{avec } \Sigma = \Fl(X) \ \ ),
\]
autant l'établir tout de suite dans le cadre général où on se donne une partie
\[
\Sigma \subseteq \Fl(X),
\]
et on définit, pour $\tau \in \Ty$,
\[
\Ch_{\Sigma,\tau}(X) \subseteq \Ch_{\tau}(X)
\]
comme la sous-catégorie pleine de $\Ch_{\tau}(X)$ formée des chemins dans $X$ dont \textit{toutes les flèches rétrogrades}

\msp{page 66}

sont dans $X$ \texttt{[plutôt dans $\Sigma$]}. On notera que pour qu'un chemin $c$ et son opposé $c^o$ soient tous les deux des $\Sigma$-\textit{chemins,} il faut et il suffit que \textit{toutes} les flèches du chemin $c$ soient dans $\Sigma$ (on dira alors que $c$ est un $\Sigma$-\textit{chemin stricte}), puisque les flèches rétrogrades de $c^o$ sont les flèches directes de $c$.

Pour que les sous-catégories pleines $\Ch_{\tau,\Sigma}(X)$ soient stables par les foncteurs images inverses
\[
f^* : \Ch_{\tau'}(X) \longrightarrow \Ch_{\tau}(X)
\]
associés aux flèches
\[
f : \tau \longrightarrow \tau'
\]
dans $\uTy$, ou simplement pour qu'y ait stabilité pour les deux foncteurs $f^*$ associés à un $f : \Delta_- \longrightarrow 1$ et $\Delta_- \longrightarrow \Delta_-^2$, il faut et il suffit que $\Sigma$ soit stable par composition et contienne toutes les flèches identiques % [sic] qu'y ait -> qu'il y ait, missing word
\[
\left\{
\begin{aligned}
& x \in \mathrm{Ob}\,X \ \Longrightarrow \ \id_x \in \Sigma \\
& u, v \in \Sigma, \ \ u,\, v \text{ composables} \Longrightarrow uv \in \Sigma.
\end{aligned}
\right.
\]
En effet, les flèches dans un chemin $f^*(c') \in \Ch_{\tau}(X)$ sont soit des flèches identiques (qui correspondent aux flèches de transition de $I_\tau$ dégénérées pour $f$), soit des flèches parmi celles de $\tau$, ou des composés d'une suite de flèches consécutives composables de $\tau'$ (qui correspondent aux flèches de transition dans $I_{\tau'}$ incluses dans la partie d'une flèche de transition de $I_\tau$), avec conservation des signatures.

\printedpage{45}
% head: A. Grothendieck : Dérivateurs, Ch. VII, Édition des Universités de Montpellier II et Paris VII
% foot: Version 22 décembre 2001

\msp{page 67}

On suppose ces conditions satisfaites, de sorte que les $\Ch_{\tau,\Sigma}(X)$ forment les fibres d'une sous-catégorie fibrée pleine de $\Ch(X)$ sur $\uTy$, qu'on dénotera par
\[
\Ch_{\Sigma}(X) \hookrightarrow \Ch(X).
\]
Il est évident que cette sous-catégorie est stable par l'opération de composition des chemins (mais non pas, en général, on l'a dit, par passage de $c$ à $c^o$ (il faudrait pour cela que $\Sigma$ contienne \textit{toutes} les flèches de $X$, puisque pour toute $u : a \longrightarrow b$, qui est donc dans $\Ch_{\Delta_+,\Sigma}(X)$, $u^o$ donné par $b \xleftarrow{\ u\ } a$ dans $\Ch_{\Delta_-}(X)$ devrait être un $\Sigma$-chemin, d'où $u \in \Sigma$)).

On va définir une catégorie canonique
\[
\Pi_\Sigma(X),
\]
ayant même objets que $X$, et un foncteur canonique
\[
X \longrightarrow \Pi_\Sigma(X)
\]
qui est l'identité sur les objets, et qui transforme les $u \in \Sigma$ en flèches inversibles. Pour cela on calque la description du groupoïde fondamental (mais $\Pi_\Sigma(X)$ ne sera pas en général un groupoïde), en posant
\[
\Hom_\Sigma(x,y) = \pi_0(\Ch_\Sigma(x,y)),
\]

\msp{page 68}

et en définissant la composition des `$\Sigma$-flèches' en termes de la composition des chemins. L'argument qui prouvait, dans le cas de $\Pi(X)$, que \textit{toute} flèche de $X$, vue comme chemin $c$ de type $\Delta_+$, était inversible dans $\Pi(X)$, ayant comme inverse $c^o$ de type $\Delta_-$, est encore valable, pourvu que $c^o$ soit encore bien un $\Sigma$-chemin, ce qui exige que $u \in \Sigma$.

On va l'interpréter en termes d'un foncteur canonique
\[
\kappa : X \longrightarrow \Pi_\Sigma(X) \qquad \text{(\footnotemark[19])}
\tag*{$(*)$}
\]
\footnotetext[19]{Sur les flèches, $\kappa$ est défini, pour $u : a \longrightarrow b$, par $\kappa(u) =$ classe d'homotopie de $\ch(u)$ dans $\Ch_\Sigma(a,b)$, où $\ch(u)$ est le chemin de type $\Delta_+$ de $a$ vers $b$ défini par $u$.}
qui se décrit encore comme dans le cas du groupoïde fondamental, en associant à toute flèche $u$ de $X$ la classe d'homotopie du chemin $c(u)$ qu'elle définit. On voit donc que ce foncteur canonique transforme les $u \in \Sigma$ en flèches inversibles.

Quand les flèches de $\Sigma$ sont déjà inversibles dans $X$, je dis que le foncteur $(*)$ est un isomorphisme, i.e. que pour tout $a, b \in \mathrm{Ob}\,X$,
\[
\alpha : \Hom_X(a,b) \longrightarrow \pi_0\Ch_\Sigma(a,b), \qquad u \longmapsto \underbrace{\ch_\Sigma(u)}_{= \ch(u)}
\]
est bijectif. Pour le voir, on va définir
\[
\beta : \pi_0(\Ch_\Sigma(a,b)) \longrightarrow \Hom_X(a,b),
\]
qui sera inverse de $\beta$. On définit % [sic] inverse de $\beta$ -> inverse de $\alpha$, $\beta$ is meant to be the inverse of $\alpha$

\printedpage{46}
% head: A. Grothendieck : Dérivateurs, Ch. VII, Édition des Universités de Montpellier II et Paris VII
% foot: Version 22 décembre 2001

\msp{page 69}

$\beta(c)$, pour une $\Sigma$-flèche de type $\tau$, le longueur $n$, en inversant toutes les flèches rétrogrades (ce qui est possibles, puisqu'elles sont dans $\Sigma$), de façon à obtenir un chemin $c^\natural$ de type $\Delta_+^n$, et en prenant le composé (dans $X$) des flèches de ce chemin $c^\natural$. Il est évident alors que $\beta\alpha = \id$, il faut prouver que l'on a aussi $\alpha\beta = \id$, i.e. que dans $\Pi_\Sigma(X)$, les flèches $u = \int c^\natural$ et $c$ sont égales, i.e. $u$ et $c$ sont homotopes dans $\Ch_\Sigma(a,b)$. Considérons le chemin $c^\natural$, de type $\Delta_+^n$ : % [sic] le longueur -> la longueur, gender; possibles -> possible, number
\[
a = a_0 \xrightarrow{u_1} a_1 \xrightarrow{u_2} a_2 \longrightarrow \ \cdots \ \longrightarrow a_{n-1} \xrightarrow{u_n} a_n = b,
\]
et soit $u_i$ ($i \in [1,n]$) une flèche de ce chemin, correspondant à une flèche de transition $\underline{u}_i$ \textit{rétrograde} de $\tau$, donc par construction
\[
u_i = v_i^{-1} \qquad \text{dans } X\text{, avec } v_i \in \Sigma
\]
la flèche d'indice $i$ dans le chemin initial $c$ de type $\tau$. Si $\underline{u}_i$ par contre est directe, on aura
\[
u_i = v_i.
\]
(NB On désigne par $v_i$ la $i$-flèche dans le chemin $c$.) Or on a
\[
c = c_n * c_{n-1} * \cdots * c_1 \qquad \text{dans } \Pi_\Sigma(X)
\tag*{$(*)$}
\]
et le facteur $c_i$ est soit la classe d'homotopie de $\ch(v_i) = \ch(u_i)$ (cas où la flèche de transition $\underline{u}_i$ de $\tau$ est directe), soit celle du chemin \textit{opposé} au chemin $\ch(v_i) : a_i \xrightarrow{v_i} a_{i-1}$ de type $\Delta_+$, donc le chemin $\ch(v_i)^o : a_{i-1} \xleftarrow{v_i} a_i$ de type $\Delta_-$. Je dis que
\[
\boxed{\ \ch(v_i)^o = \ch(u_i) \qquad \text{\textit{dans} } \Pi(X),\ }
\tag*{(1)}
\]
ce qui prouve que dans $\Pi_\Sigma(X)$, $c$ est égal au composé des $\ch(u_i)$ ($i \in [1,n]$), donc à $c^\natural$. Il faut prouver de plus que
\[
c^\natural = \ch(u) \qquad (u = u_nu_{n-1}\cdots u_1) \qquad \text{dans } \Pi(X).
\tag*{(2)}
\]

\msp{page 70}

Donc il y a deux vérifications à faire (où je suppose à nouveau $\Sigma$ quelconque) :

\begin{enumerate}[label={}]
\item[$1^o$)] Soit $u \in \Sigma$, $u : a \longrightarrow b$, supposons $\Sigma$ inversible dans $X$. Alors
\[
\ch(u)^o = \ch(u^{-1}) \qquad \text{dans } \Pi_\Sigma(X),
\]
i.e. les deux chemins de $b$ vers
\[
\begin{array}{cccl}
b & \xleftarrow{\ u\ } & a & \text{de type } \Delta_- \\
b & \xrightarrow{u^{-1}} & a & \text{de type } \Delta_+
\end{array}
\]
\printedpage{47}
% head: A. Grothendieck : Dérivateurs, Ch. VII, Édition des Universités de Montpellier II et Paris VII
% foot: Version 22 décembre 2001
sont homotopes dans $\Ch_\Sigma(b,a)$. Or on sait déjà, comme $\kappa : X \longrightarrow \Pi_\Sigma(X)$ est un foncteur, que
\[
\hot\,\ch(u^{-1}) = (\hot\,\ch(u))^{-1} \qquad \text{dans } \Pi_\Sigma(X)
\]
où $\hot(\bullet)$ désigne la classe d'homotopie d'une chaine (ici classe dans $\Ch_\Sigma(b,a)$ et dans $\Ch_\Sigma(a,b)$ respectivement). Donc l'énoncé à prouver c'est que $\hot(\ch(u)^o)$ est égal à l'\textit{inverse} de $\hot(\ch(u))$ dans $\Pi_\Sigma(X)$. C'est donc ce qu'on a prouvé pour \textit{tout} $u \in \Fl(X)$, dans le cas de $\Pi(X)$. L'argument doit marcher, du moment qu'il a un sens, i.e. que $\ch(u)^o$ est bel et bien un $\Sigma$-chemin, i.e. $u \in \Sigma$.

\item[$2^o$)] La formule (2) signifie que pour tout chemin $c$ de type $\Delta_+^n$ (donc tautologiquement un \msp{page 71} $\Sigma$-chemin), allant de $a$ vers $b$, sa classe dans $\Hom_\Sigma(a,b)$ n'est autre que celle de $\ch(u)$, où $u$ est le composé $u_nu_{n-1}\cdots u_1$. Mais on a
\[
c = c_1 * c_2 * \cdots * c_n,
\]
d'où dans $\Pi_\Sigma(X)$
\[
\hot(c) = \hot(c_n)\hot(c_{n-1})\cdots\hot(c_1),
\]
or $\hot(c_i) = \hot\,\ch(u_i) = \kappa(u_i)$, donc
\[
\begin{aligned}
\hot(c) \ &= \ \kappa(u_n)\kappa(u_{n-1})\cdots\kappa(u_1) \\
&= \ \kappa(\underbrace{u_nu_{n-1}\cdots u_1}_{u}),
\end{aligned}
\]
puisque $\kappa$ est un foncteur, donc
\[
\hot(c) = \hot\,\ch(u),
\]

\hfill c.q.f.d.
\end{enumerate}

Notons maintenant que si $(X', \Sigma')$ est un deuxième couple d'une catégorie $X'$ et d'un ensemble \textit{`admissible'} de flèches de $X'$, et si
\[
f : X \longrightarrow X'
\]
est un foncteur tel que
\[
f(\Sigma) \subseteq \Sigma',
\]
alors il existe un \textit{unique} foncteur
\[
\bar{f} = \Pi_\Sigma(f) : \Pi_\Sigma(X) \longrightarrow \Pi_{\Sigma'}(X')
\]
qui `prolonge' $f$, i.e. tel que le carré
\printedpage{48}
% head: A. Grothendieck : Dérivateurs, Ch. VII, Édition des Universités de Montpellier II et Paris VII
% foot: Version 22 décembre 2001
\msp{page 72}
\[
\begin{tikzcd}
X \arrow[r, "\kappa_X"] \arrow[d, "{\scriptstyle f}"] & \Pi_\Sigma(X) \arrow[d, "{\scriptstyle \bar{f} = \Pi_\Sigma(f)}"] \\
X' \arrow[r, "\kappa_{X'}"] & \Pi_{\Sigma'}(X')
\end{tikzcd}
\]
soit commutatif. L'unicité résulte du fait que sur les objets $\bar{f}$ doit être donné par $f$, et que sur les flèche la commutativité du carré implique que $\bar{f}$ est connu sur les flèches de la forme $\kappa_X(u)$, ainsi donc que sur les inverses de telles flèches, quand ces inverses sont inversibles, donc en particulier sur les $\kappa_X(u)^{-1}$ pour $u \in \Sigma$. Or toute flèche de $\Pi_\Sigma(X)$, provenant d'un chemin, est composé d'une suite de flèches de cette forme. Pour l'existence de $\bar{f}$, on note que $f$ induit des foncteurs % [sic] les flèche -> les flèches, plural
\[
\Ch_\tau(X; a, b) \longrightarrow \Ch_\tau(X'; f(a), f(b))
\]
qui commutent aux opérations $\varphi^*$ sur les chemins associées aux foncteurs
\[
\varphi : \tau \longrightarrow \tau',
\]
d'où un homomorphisme de catégories fibrées sur $\uTy$
\[
f_{a,b} : \Ch(X; a, b) \longrightarrow \Ch(X'; f(a), f(b)).
\]

\msp{page 73}

D'où une application sur les $\pi_0$, i.e.
\[
\bar{f}_{a,b} : \Hom_\Sigma(a,b) \longrightarrow \Hom_\Sigma(f(a), f(b)),
\]
donnée par
\[
\bar{f}_{a,b}(\hot(c)) = \hot(f_{a,b}(c)).
\]
On a évidemment, pour deux chemins $\gamma : a \longrightarrow b$ et $\gamma' : b \longrightarrow c$,
\[
f_{a,c}(\gamma\gamma') = f_{b,c}(\gamma')f_{a,b}(\gamma),
\]
ce qui prouve que $\bar{f}$ est compatible avec la composition. Il est trivial qu'il est compatible avec les identités (car $f_{a,a}$ transforme le chemin vide en $a$ en le chemin vide en $f(a)$), donc on a un foncteur $\bar{f}$, et il rend bel et bien le carré commutatif.

On peut maintenant prouver la propriété universelle de $\Pi_\Sigma(X)$ et de
\[
X \xrightarrow{\ \kappa\ } \Pi_\Sigma(X),
\]
savoir que tout foncteur $f : X \longrightarrow X'$ tel que $u \in \Sigma \Longrightarrow f(u)$ inversible, se factorise \textit{de façon unique} par $\Pi_\Sigma(X)$. L'unicité se voit comme tantôt (\footnote[20]{C'est inutile, cf. ci-dessous.}). Pour l'existence, on note que $f$ applique $\Sigma$ dans l'ensemble $\Sigma'$

% ---- from Der-ChVII-456_p051-055.tex
\setcounter{footnote}{20}

\printedpage{49}
% head: A. Grothendieck : Dérivateurs, Ch. VII, Édition des Universités de Montpellier II et Paris VII
\msp{page 74}
des flèches inversibles de $X'$ (lequel ensemble satisfait bien aux deux conditions d'admissibilité). Donc par ce qui précède on a un diagramme commutatif
\[
\begin{tikzcd}
X \ar[r,"\kappa"] \ar[d,"f"'] & \Pi_{\Sigma}(X) \ar[d,"\bar f = \Pi_{\Sigma}(f)"] \ar[dl,dotted,"\tilde f"'] \\
X' \ar[r,"\kappa'","\sim"'] & \Pi_{\Sigma'}(X'),
\end{tikzcd}
\]
où $\bar f$ est uniquement déterminé. Comme $\kappa'$ est un isomorphisme (comme on a vu tantôt, puisque $\Sigma'$ \texttt{[est]} formé des flèches inversibles), l'existence et l'unicité de $\bar f$ équivaut à la conclusion annoncée,

\hfill c.q.f.d.

\textbf{Encore un sorite} sur les classes d'homotopie d'éléments d'une catégorie $X$. De façon générale, on voit que $X$ et $\Pi_{\Sigma}(X) = X\Sigma^{-1}$ ont mêmes composantes connexes, car si
\msp{page 75}
$a$ et $b$ dans $X$ sont reliés par une flèche de $X\Sigma^{-1}$, i.e. par un $\Sigma$-chemin, alors ils sont dans une même composante connexe de $X$. Mais l'inverse est également vrai, comme on voit p.~ex. en prenant $\Sigma = \mathrm{Fl}(X)$, et en disant que dans le \textit{groupoïde} $\Pi(X)$, où deux sommets sont dans une même composante \texttt{[connexe]} si et seulement si $\exists\ x \longrightarrow y$ (relation qui est non seulement transitive et reflexive, mais aussi symétrique dans un groupoïde). Un chemin
\[
a \xrightarrow{\quad c \quad} b
\]
sera appelé aussi une \textit{homotopie} de $a$ vers $b$. On dit que $a$ et $b$ sont \textit{homotopes} s'il existe une homotopie $c$ de $a$ vers $b$, auquel cas il existe aussi une homotopie `inverse' $c^{o}$ de $b$ vers $a$. Cela signifie aussi que $a$ et $b$ sont dans la même composante connexe de $X$, i.e. définissent le même élément de $\pi_0(X)$.

L'homotopie sera dite \textit{simple} si $c$ est de type $\Delta_{+}$ ou $\Delta_{-}$, i.e. si elle est donnée par une flèche $u : a \longrightarrow b$ ou $v : b \longrightarrow a$. Si on cherche une homotopie entre $a$ et $b$, on essaiera de la trouver la plus `courte' possible. Si on a une homotopie $c$ de longueur $l$ de $a$ vers $b$, alors en remplaçant $\tau$ par $\tau_{\red}$, donc $I_{\tau}$ par $I_{\tau\,\red}$, le sous-intervalle réduit
\msp{page 76}
correspondant, on trouve une homotopie de $a$ à $b$ de type $\tau_{\red}$, dont la longueur est égale à la \textit{`longueur utile'} du type $\tau$ de $c$ (ce qui justifie la terminologie). Notons que si
\[
\tau = \Delta_{\varepsilon}^{l_1}\Delta_{-\varepsilon}^{l_2}\Delta_{\varepsilon}^{l_3}\cdots\Delta_{(-1)^{r}\varepsilon}^{l_r},
\]
alors
\[
\tau_{\red} = \Delta_{\varepsilon}\Delta_{-\varepsilon}\Delta_{\varepsilon}\cdots\Delta_{(-1)^{r}\varepsilon},
\]
% foot: Version 22 décembre 2001
\printedpage{50}
% head: A. Grothendieck : Dérivateurs, Ch. VII, Édition des Universités de Montpellier II et Paris VII
\texttt{[plutôt deux fois} $(-1)^{r-1}$\texttt{]} où $r$ est la longueur utile de $\tau$ (la longueur pas utile étant $n = \sum\limits_{i\in[1,r]} l_i$).

Quand on part d'une homotopie entre \textit{foncteurs}
\[
X \longrightarrow Y,
\]
il est entendu qu'il s'agit d'homotopies dans la catégorie $\uHom(X,Y)$. (Ou, suivant le contexte, dans une sous-catégorie précisée de $\uHom(X,Y)$.)

\msp{page 77}
À la page 61, j'ai invoqué l'existence d'un \textit{foncteur} de \textit{`composition des chemins'}
\[
\Ch_{\Sigma}(X;a,b)\times\Ch_{\Sigma}(X;b,c) \ \longrightarrow\ \Ch_{\Sigma}(X;a,c),
\]
qui sur la fibre du premier membre au dessus de $(\tau,\tau') \in \operatorname{Ob}(\Ty\times\Ty)$, serait donné par
\[
\Ch_{\Sigma,\tau}(X;a,b)\times\Ch_{\Sigma,\tau'}(X;b,c) \ \longrightarrow\ \Ch_{\Sigma,\tau\tau'}(X;a,c), \tag*{$(*)$}
\]
la composition des chemins de types précisés $\tau$, $\tau'$. Qu'on ait là un \textit{foncteur,} pas seulement une application sur les ensembles d'objets, est assez évident, et peut s'exprimer par exemple à l'aide du diagramme d'amalga\texttt{[ma]}tion
\[
\begin{tikzcd}[column sep=small]
 & \varepsilon = \{0,1\} \ar[dl,hook'] \ar[dr,hook] & \\
I_{\tau} \ar[dr,hook] & & I_{\tau'} \ar[dl,hook'] \\
 & I_{\tau\tau'} &
\end{tikzcd}
\]
en le transformant par $\uHom(\cdot,X)$ (au lieu de le transformer par $\mathrm{Hom}_{\mathrm{Cat}}(\cdot,X)$ seulement), du moins dans le cas particulier où $\Sigma = \mathrm{Fl}(X)$ (qui était celui de la page 61). Mais le passage aux flèches est assez évident - pour deux composés $c'_1\circ c_1$ et $c'_2\circ c_2$, de
\[
\begin{tikzcd}[column sep=small]
a \ar[rr,"c_1"] \ar[d,equal] & {} \ar[d,"w"'] & b \ar[rr,"c'_1"] \ar[d,equal] & {} \ar[d,"w'"] & c \ar[d,equal] \\
a \ar[rr,"c_2"] & {} & b \ar[rr,"c'_2"] & {} & c
\end{tikzcd}
\]
\texttt{[on]} voit que les homomorphismes $w''$ du premier dans le second, dans $\Ch_{\Sigma}(X;a,c)$, qui sont l'identité sur le sommet d'indice $n+1$ des deux chemins (où $n = \operatorname{long}\tau$), correspondant à $1_{\tau}$ et à $0_{\tau'}$, `sont' exactement les couples
% foot: Version 22 décembre 2001
\printedpage{51}
% head: A. Grothendieck : Dérivateurs, Ch. VII, Édition des Universités de Montpellier II et Paris VII
\msp{page 78}
formés d'un homomorphisme $w : c_1 \longrightarrow c_2$ et $w' : c'_1 \longrightarrow c'_2$ dans $\Ch_{\Sigma,\tau}(a,b)$ et $\Ch_{\Sigma,\tau'}(b,c)$ respectivement, et si $w''$, $w$, $w'$ se correspondent ainsi, on posera
\[
w'' = w * w'.
\]
Il y a une transitivité évidente à vérifier pour voir que ça définit un \textit{foncteur} $(*)$. On peut exprimer la situation à l'aide du diagramme
\[
\begin{tikzcd}[column sep=small]
 & \Ty\times\Ty \ar[rrr,"\pi\ (\text{composition})"] \ar[dl,"p_1"'] \ar[dr,"p_2"] & & & \Ty \\
\Ty & & \Ty, & &
\end{tikzcd}
\]
où $p_1$, $p_2$ sont les deux projections. On considère sur $\Ty\times\Ty$ les trois catégories fibrés % [sic] fibrés -> fibrées, agreement with "catégories"
\[
p_1^*(\CH_{\Sigma,\tau}(a,b)),\ \ p_2^*(\CH_{\Sigma,\tau'}(b,c)),\ \ \pi^*(\Ch_{\Sigma,\tau\tau'}(a,c))
\]
\texttt{[plutôt} $\pi^*(\CH_{\Sigma,\tau\tau'}(a,c)$\texttt{]}, dont les fibres en $(\tau,\tau')$ sont respectivement % [sic] \pi^*(\CH_{\Sigma,\tau\tau'}(a,c) -> \pi^*(\CH_{\Sigma,\tau\tau'}(a,c)), one closing parenthesis missing
\[
\Ch_{\Sigma,\tau}(a,b),\ \ \Ch_{\Sigma,\tau'}(b,c),\ \ \Ch_{\Sigma,\tau\tau'}(a,c).
\]
(On écrit $\CH$ au lieu de $\Ch$, pour indiquer que l'on considère les catégories envisagées comme des catégories (fibrées) \textit{au dessus de} $(\Ty)$ (\footnote{Mais il s'agit pour $p_1^*$, $p_2^*$ et $\pi^*$ de produits cartésiens ordinaires dans Cat, il n'y a pas à avoir peur et à écrire des $\CH$ \ldots}).) Au dessus de $\Ty\times\Ty$, on a un homomorphisme (cartésien) de catégories fibrées
\[
p_1^*(\CH_{\tau}(a,b))\times_{\Ty\times\Ty} p_2^*(\CH_{\tau}(b,c)) \ \longrightarrow\ \pi^*(\CH_{\tau}(a,c)) \tag*{$(**)$}
\]
\texttt{[plutôt} $p_1^*(\CH_{\tau}(a,b))\times_{\Ty\times\Ty} p_2^*(\CH_{\tau'}(b,c)) \ \longrightarrow\ \pi^*(\CH_{\tau\tau'}(a,c))$\texttt{]},
\msp{page 79}
car c'est OK sur les fibres ($(*)$ page 77), et si
\[
(\varphi,\varphi') : (\tau_1,\tau'_1) \longrightarrow (\tau,\tau') \qquad (\varphi : \tau_1 \longrightarrow \tau,\ \varphi' : \tau'_1 \longrightarrow \tau')
\]
est une flèche dans $\Ty\times\Ty$, on en déduit un diagramme commutatif
\[
\begin{tikzcd}
\Ch_{\Sigma,\tau}(a,b)\times\Ch_{\Sigma,\tau'}(b,c) \ar[r,"\text{comp.}"] \ar[d,"\varphi^*_{a,b}\times(\varphi'_{b,c})^*"] & \Ch_{\Sigma,\tau}(a,c) \ar[d,"(\varphi*\varphi')^*_{a,c}"] \\
\Ch_{\Sigma,\tau_1}(a,b)\times\Ch_{\Sigma,\tau'_1}(b,c) \ar[r,"\text{comp.}"] & \Ch_{\Sigma,\tau_1\tau'_1}(a,c)
\end{tikzcd}
\]
% foot: Version 22 décembre 2001
\printedpage{52}
% head: A. Grothendieck : Dérivateurs, Ch. VII, Édition des Universités de Montpellier II et Paris VII
\texttt{[plutôt} $\Ch_{\Sigma,\tau\tau'}(a,c)$\texttt{]}, i.e. la composition des chemins d'extrémités fixés est compatible avec les images inverses de chemins par flèches entre les types de chemins.

Mais l'homomorphisme cartésien $(**)$ de catégories fibrées sur $\Ty\times\Ty$ n'est pas tout à fait la même chose qu'un anodin foncteur
\[
\Ch_{\Sigma}(a,b)\times\Ch_{\Sigma}(b,c) \ \longrightarrow\ \Ch_{\Sigma}(a,c). \tag*{$(*)$}
\]
La formule $(*)$ définit en tout cas une application sur les ensembles d'objets, puisque les ensembles d'objets de chacune des trois catégories envisagées se décompose en somme d'ensembles $\Ch_{\Sigma,\tau}(\bullet,\bullet)$, où $\tau$ parcourt $\mathrm{Ty} = \operatorname{Ob}\Ty$. Mais cela se prolonge-t-il de façon naturelle aux morphismes?

\msp{page 80}
Mais voilà que je me rappelle que ce fameux produit cartésien $p_1^*(\text{fourbi})\times p_2^*(\text{machin})$ n'est autre que le produit cartésien tout bête
\[
\Ch_{\Sigma}(a,b)\times\Ch_{\Sigma}(b,c),
\]
envoyé sur $\Ty\times\Ty$ par le produit cartésien des deux foncteurs de projection
\[
\begin{aligned}
p_{a,b} : \Ch_{\Sigma}(a,b) &\ \longrightarrow\ \ \Ty\\
p_{b,c} : \Ch_{\Sigma}(b,c) &\ \longrightarrow\ \ \Ty,
\end{aligned}
\]
d'où
\[
\begin{tikzcd}
\Ch_{\Sigma}(a,b)\times\Ch_{\Sigma}(b,c) \ar[d,"p_{a,b}\times p_{b,c}"] \\
\Ty\times\Ty
\end{tikzcd}
\]
\msp{page 81}
qui est bel et bien fibrant (comme produit cartésien de deux foncteurs fibrants) et a les fibres qu'il faut! On trouve donc bien, par le sorite général sur les foncteurs fibrants (ici sur $\Ty\times\Ty$) un foncteur de catégories fibrées sur $\Ty\times\Ty$
\[
\Ch_{\Sigma}(a,b)\times\Ch_{\Sigma}(b,c) \ \longrightarrow\ \qquad \underbrace{\pi^*(\CH_{\Sigma}(a,c))}_{=\,(\Ty\times\Ty)\times_{\Ty}\Ch_{\Sigma}(a,c)\ \ \text{(\footnotemark)}} \qquad .
\]
\footnotetext{produit fibré par $\pi : \Ty\times\Ty \longrightarrow \Ty$.}
Comme on a d'autre part un foncteur du dernier membre (produit fibré) vers le facteur $\Ch_{\Sigma}(a,c)$, on gagne. (On l'aurait retrouvé par le raisonnement direct que je m'apprêtais de faire, et que j'ai biffé - mais je suis plus heureux comme ça!)

% foot: Version 22 décembre 2001
\printedpage{53}
% head: A. Grothendieck : Dérivateurs, Ch. VII, Édition des Universités de Montpellier II et Paris VII
Pour bien faire, il faudrait prouver l'\textit{associativité} des foncteurs de composition, pas seulement au niveau des objets, par des arguments comme ceci, mais ça revient simplement à une relation du type
\[
(w''*w')*w = w''*(w'*w)
\]
pour un système de trois flèches dans
\[
\Ch_{\Sigma}(a,b),\ \Ch_{\Sigma}(b,c)\ \ \text{et}\ \ \Ch_{\Sigma}(c,d)
\]
respectivement. Ça ne peut être que trivialement
\msp{page 82}
vrai, mais fastidieux à expliciter.

Les sueurs froides sont terminées!

En résumé, le formalisme des $\Ch_{\Sigma}(X)$ et des $\Ch_{\Sigma}(X;a,b)$ développé jusqu'à présent donne satisfaction complète sur les points suivants :

\begin{enumerate}
\item[a)] Bon formalisme de composition des chemins (qui était dégueulasse avant).

\item[b)] Traitement élégant de la construction des catégories de fractions $\mathcal{M}\Sigma^{-1}$. Mais je m'attends que ce sera moins joli qu'avant, quand il s'agira de remplacer un $(\mathcal{M},\Sigma)$ par des $(\ \underbrace{\mathcal{M}(X)}_{=\,\uHom(X,\mathcal{M})}\ ,\Sigma_X)$ ($X$ une catégorie d'indices) \texttt{[plutôt} $\uHom(X^{o},\mathcal{M})$\texttt{]}. C'est pour y pallier qu'il me faudra développer des énoncés de comparaison, entre la catégorie $\Ch_{\Sigma}$ construite ici, et d'autres dont le type dans Hot se prête aux calculs de façon plus directe, comme dans l'état premier de la question. De plus, on sera sûrement obligé d'en venir à des Ind-catégories, pour pouvoir suivre la variance de $\mathcal{M}(X)\Sigma_X^{-1}$ pour $X$ variable.

\item[c)] Bon formalisme des `homotopies' et des compositions des homotopies (ce qui est une autre façon de parler des chemins et de leur composition).
\end{enumerate}
% foot: Version 22 décembre 2001

% ---- from Der-ChVII-456_p056-060.tex
\printedpage{54}
% head: A. Grothendieck : Dérivateurs, Ch. VII, Édition des Universités de Montpellier II et Paris VII
% foot: Version 22 décembre 2001

\msp{page 83}

\section*{8\quad Le lemme d'homotopie pour les morphismes d'intervalles}

C'est le suivant :

\textbf{Lemme 1.} \textit{Soient $\tau,\tau'\in\mathrm{Ty}$, et $f,g:\tau\rightrightarrows\tau'$ deux morphismes dans $\uTy$. Alors $f$, $g$ sont homotopes dans la catégorie $\uHom\,\underline{\mathrm{int}}(I_\tau,I_{\tau'})$ des morphismes croissants d'intervalles orientés de $I_\tau$ dans $I_{\tau'}$.}

Cette catégorie pourrait licitement se désigner par $\uHom_{\uTy}(\tau,\tau')$, ou $\uHom(\tau,\tau')$. C'est une sous-catégorie de la catégorie $\uHom(I_\tau,I_{\tau'})$ de \textit{tous} les foncteurs de $I_\tau$ dans $I_{\tau'}$, et une sous-catégorie pleine de la catégorie $\uHom^\varepsilon(I_\tau,I_{\tau'})$ des foncteurs $f$ qui sont compatibles avec la donnée des sources et des buts, i.e. qui rendent commutatif le diagramme
\[
\begin{tikzcd}[column sep=small]
 & \varepsilon=\{0,1\} \arrow[dl, hook', "{(\sigma_\tau,\beta_\tau)}"'] \arrow[dr, hook, "{(\sigma_{\tau'},\beta_{\tau'})}"] & \\
I_\tau \arrow[rr, "f"] & & I_{\tau'}.
\end{tikzcd}
\]

Si $f$, $g$ sont dans cette catégorie (pas forcément morphismes croissants d'ailleurs), une flèche de $f$ dans $g$ est donc la donnée, pour tout $i\in I_\tau$, d'une flèche $f(i)\xrightarrow{w_i}g(i)$ dans $I_{\tau'}$, avec les deux conditions :

\begin{enumerate}[label={}, leftmargin=*]
\item[a)] Si $i\in\{0_\tau,1_\tau\}$, $w_i=\mathrm{id}$, et
\item[b)] le carré commutatif habituel pour une flèche $i\longrightarrow j$ dans $I_\tau$. (Il suffit d'avoir cette compatibilité pour les flèches de transition $\ul{u}_i:i-1\longrightarrow i$ dans $I_\tau$ ($i\in[1,\operatorname{long}\tau]$).)
\end{enumerate}

Une autre façon de décrire cette catégorie $\uHom^\varepsilon$, c'est de dire que c'est la catégorie
\msp{page 84}
\[
\underline{\mathrm{Ch}}_\tau(I_{\tau'};0_{\tau'},1_{\tau'})\qquad(\simeq\uHom^\varepsilon(I_\tau,I_{\tau'})),
\]
mais $\uHom_{\uTy}(\tau,\tau')$ est la sous-catégorie pleine de cette catégorie formée des \textit{chemins croissants} dans $I_{\tau'}$. Le lemme d'homotopie signifie aussi que cette catégorie $\uHom_{\uTy}(\tau,\tau')$ est toujours \textit{connexe.} Elle n'est d'ailleurs pas forcément 0-connexe, car elle peut être vide. Notons ici, pour plus tard, la

\textbf{Proposition 1 :} \textit{Soient $\tau,\tau'\in\mathrm{Ty}$, soient $r$, $r'$ leurs longueurs utiles. Pour que $\mathrm{Hom}_{\uTy}(\tau,\tau')$ soit non vide, il faut et il suffit que l'une ou l'autre des deux conditions suivantes soit satisfaite :}

\begin{enumerate}[label={}, leftmargin=*]
\item[\textit{a)}] $r\geq r'+1$

\printedpage{55}
% head: A. Grothendieck : Dérivateurs, Ch. VII, Édition des Universités de Montpellier II et Paris VII
% foot: Version 22 décembre 2001

\item[\textit{b)}] \textit{$r=r'$, et (si $r\neq0$) $\tau$ et $\tau'$ ont même direction initiale (de sorte que $\tau_{\red}=\tau'_{\red}$).}
\end{enumerate}

(Démonstration $\pm$ immédiate.)

Notons que si $\tau$, $\tau'$ n'ont pas même direction initiale, alors tout morphisme $\tau\longrightarrow\tau'$ est forcément dégénéré sur le premier intervalle critique de $I_\tau$. On aura donc $\tau=\Delta_\varepsilon^{l_0}\cdot\tau_0$ (où $l_0$ est la longueur du premier intervalle critique, et où $\tau_0$ a \texttt{[la]} direction initiale opposée de celle de $\tau$), et un isomorphisme
\[
\uHom(\tau,\tau')\simeq\uHom(\tau_0,\tau'),
\]
où cette fois $\tau_0$, $\tau'$ ont même direction initiale. On supposera donc désormais que \textit{la direction initiale de $\tau$, $\tau'$ est la même.} La condition $\uHom(\tau,\tau')\neq\emptyset$ équivaut maintenant à $\underline{r\geq r'}$, ce que nous pouvons supposer. Je vais alors décrire un élément canonique
\msp{page 85}
$F_0$ de $\uHom(\tau,\tau')$, ainsi. Soient
\[
I_1,I_2,\dots,I_r
\]
les intervalles critiques de $I_\tau$, et
\[
I'_1,I'_2,\dots,I'_{r'}\qquad(r'\leq r)
\]
ceux de $I_{\tau'}$. Le foncteur $F_0$ proviendra de morphismes d'intervalles orientés (comme
\[
\left\{
\begin{array}{ccc}
I_1 & \xrightarrow{f_1} & I'_1\\[4pt]
I_2 & \xrightarrow{f_2} & I'_2\\[4pt]
 & \dots & \\[4pt]
I_{r'} & \xrightarrow{f_{r'}} & I'_{r'}
\end{array}
\right.\qquad)
\]
et sera donc forcément trivial sur les intervalles critiques $I_\alpha$ pour $r'<\alpha\leq r$ (il y en a si et seulement si $r>r'$). Comme $I_\tau$, $I_{\tau'}$ ont même direction initiale, il s'ensuit que pour $\alpha\in[1,r']$, $I_\alpha$ et $I'_\alpha$ ont même signature (à savoir $\varepsilon_\alpha=(-1)^{\alpha+1}\varepsilon_1$, où $\varepsilon_1$ est la signature commune de $I_1$ et $I'_1$). Donc
\[
I_\alpha=\Delta^{l_\alpha}_{\varepsilon_\alpha},\qquad I'_\alpha=\Delta^{l'_\alpha}_{\varepsilon_\alpha},
\]
et on définit $f_\alpha$ comme
\[
f_\alpha=\mathfrak{f}_{\varepsilon_\alpha,l'_\alpha}*\triv_\alpha,
\]
où
\[
\mathfrak{f}_{\varepsilon_\alpha,l'_\alpha}:\Delta_{\varepsilon_\alpha}\longrightarrow\Delta^{l'_\alpha}_{\varepsilon_\alpha}
\]
et où
\[
\triv_\alpha:\Delta^{l_{\alpha-1}}_{\varepsilon_\alpha}\longrightarrow1. % [sic] l_{\alpha-1} -> l_\alpha-1, the exponent is printed as l with subscript \alpha-1, while I_\alpha = \Delta^{l_\alpha} = \Delta * \Delta^{l_\alpha-1}
\]
En d'autres termes, pour $\varepsilon=+1$ (p. ex.) on doit définir un morphisme des intervalles type $\underbrace{\mathbf{\Delta}^+_{l_\alpha}}_{=[0,l_\alpha]}\longrightarrow\underbrace{\mathbf{\Delta}^+_{l'_\alpha}}_{=[0,l'_\alpha]}$, et on pose (avec les numérotations internes à $\mathbf{\Delta}_{l_\alpha}$, $\mathbf{\Delta}_{l'_\alpha}$)
\[
f_\alpha(0)=0,\qquad f_\alpha(i)=l'_\alpha\ \text{ pour }i\in[1,l_\alpha].
\]

\printedpage{56}
% head: A. Grothendieck : Dérivateurs, Ch. VII, Édition des Universités de Montpellier II et Paris VII
% foot: Version 22 décembre 2001

\msp{page 86}
L'idée donc, c'est d'être dans les dispositions d'un voyageur qui, tenu à un plan de voyage par étapes décrit par $\tau$, voudrait arriver aussi vite que possible à $1_{\tau'}$, en partant de $0_{\tau'}$. Il va dès la première étape $\Delta_\varepsilon\subseteq I_1$, se porter jusqu'au bout de $I'_1$, regrettant ne pas pouvoir aller plus loin (à cause du plan de voyage). Il doit rester inactif jusqu'à ce qu'il soit arrivé au moment marqué par l'extrémité finale de $I_1$ (si celle-ci est au delà de celle de sa première flèche de transition, i.e. si $l_1>1$), puis il fait son deuxième grand pas jusqu'à l'extrémité terminale de $I'_2$ etc. Il peut espérer, à chaque moment de son voyage, être en avance (ou du moins pas en retard) sur les autres voyageurs, qui auraient suivi de tactiques de progression différentes. Si cet espoir était justifié, cela signifierait que $F_0$ serait un plus grand élément de $\uHom(\tau,\tau')$ (\textbf{NB} toutes ces catégories $\uHom$ entre intervalles sont ‘ordonnées’, i.e. il y a au plus une flèche entre deux objets), i.e. un objet final. Cela impliquerait que $\uHom(\tau,\tau')$ serait même contractile - chose peut-être vraie, mais sûrement pas de façon si simpliste je crains. Je me suis d'ailleurs canulé : $F_0$ est sans doute bien un ‘plus grand’ élément, mais pour l'ordre défini par la numérotation des sommets de $I_{\tau'}$ (donc $F\leq G$ si et seulement si $F(i)\leq G(i)$ $\forall i\in I_\tau$), et nullement pour l'ordre défini par la structure de catégorie de $\uHom(\tau,\tau')$.

\msp{page 87}
L'idée pour prouver la connexité de $\uHom(\tau,\tau')$, c'est de ramener tout $F\in\mathrm{Hom}(\tau,\tau')$, par une suite d'homotopies simples, à la forme normale. Rappelons que $F$ est donné par une suite croissante de flèches de transition
\[
\ul{u}_{i_1},\ \ul{u}_{i_2},\ \dots\ ,\ \ul{u}_{i_s},\qquad i_1<i_2<\cdots<i_s,
\]
et une suite d'entiers $\geq1$
\[
\lambda_1,\ \lambda_2,\ \dots\ ,\ \lambda_s
\]
avec la condition que
\[
\Delta^{\lambda_1}_{\varepsilon_{i_1}}\Delta^{\lambda_2}_{\varepsilon_{i_2}}\cdots\Delta^{\lambda_s}_{\varepsilon_{i_s}}=\tau'. \tag*{$(*)$}
\]
Les $\ul{u}_{i_\alpha}$ sont les flèches de transition de $I_\tau$ \textit{critiques} pour $F$, i.e. non dégénérées pour $F$, et $\lambda_\alpha$ est la \textit{partie} de $\ul{u}_{i_\alpha}$ pour $F$. Le foncteur normalisé (‘pressé’) $F_0$ est distingué par les conditions suivantes a) et b) \texttt{[et c)]}. Tout d'abord :

\begin{enumerate}[label={}, leftmargin=*]
\item[a)] La suite des signes $\varepsilon_{i_1},\varepsilon_{i_2},\dots,\varepsilon_{i_s}$ est ‘alternée’, i.e. $\varepsilon_{i_\alpha}\neq\varepsilon_{i_{\alpha+1}}$ ($\alpha\in[1,s-1]$). (Cela traduit le fait que le voyageur pressé ne fait pas en plusieurs pas ce qu'il pourrait faire en un seul.)

Comme le premier de ces signes est forcément $\varepsilon'_1=\varepsilon_1$, cela signifie que la suite des signes est la même que celle des $\varepsilon_1$, $\varepsilon_2$, \dots, $\varepsilon_{r'}$ (à la longueur près). Mais mieux, que l'expression $(*)$ de $\tau'$ est son unique expression réduite, donc que $s=r'$ (longueur utile de $\tau'$), et que les $\lambda_i$ sont les longueurs $l'_i$ des intervalles critiques de $I_{\tau'}$.

\printedpage{57}
% head: A. Grothendieck : Dérivateurs, Ch. VII, Édition des Universités de Montpellier II et Paris VII
% foot: Version 22 décembre 2001

\msp{page 88}
Soit donc $\alpha\in[1,s=r']$, et considérons la flèche de transition $F$-critique $\ul{u}_{i_\alpha}$, soit $j_\alpha$ l'indice marquant l'intervalle critique de $I_\tau$ contenant la partie de $\ul{u}_{i_\alpha}$ dans $I_\tau$,
\[
\ul{u}_{i_\alpha}\in\mathrm{Fl}(I_{j_\alpha}),
\]
alors $F$ induit des morphismes (croissants) d'intervalles orientés
\[
F_\alpha=F|I_{j_\alpha}:I_{j_\alpha}\longrightarrow I'_\alpha
\]
et est trivial sur les autres intervalles critiques (ceux pas de la forme $I_{j_\alpha}$). La donnée des $F_\alpha$ détermine donc $F$ sur les $I_{j_\alpha}$, donc sur \textit{tous} les $I_j$, donc partout (par l'opération $*$ de composition par recollement). De plus la condition a) assure que chaque $F_\alpha$ est de nature assez particulière : une seule flèche de transition critique $\ul{u}_{i_\alpha}$. De plus, pour qu'on ait $F=F_0$, il faudra :

\item[b)] $\forall\alpha\in[1,r']$, $\ul{u}_{i_\alpha}$ est la \textit{première} flèche de transition de l'intervalle critique $I_{j_\alpha}$ qui la contient.

Et enfin :

\item[c)] La suite des intervalles critiques de $I_\tau$
\[
I_{j_1},\ I_{j_2},\ \dots\ ,\ I_{j_{r'}}\qquad(\text{de signes alternés }\varepsilon_1,\ -\varepsilon_1,\ \varepsilon_1,\ \ldots)
\]
associés aux flèches $F$-critiques, n'est autre que celle
\[
I_1,\ I_2,\ \dots\ ,\ I_{r'}
\]
des $r'$ premiers intervalles critiques de $I_\tau$, i.e.
\[
(j_1,j_2,\dots,j_{r'})=(1,2,\dots,r').
\]
\end{enumerate}

\msp{page 89}
Ceci vu, l'idée naturelle pour ‘déformer’ $F$ en $F_0$, c'est de trouver une suite d'homotopies qui transforment $F$ de façon à lui faire acquérir successivement les caractéristiques a), b), c) du voyageur pressé. Au besoin, pour aller plus vite au but, on utilisera l'hypothèse de récurrence sur $n=\operatorname{long}\tau$, que le théorème de connexité est déjà prouvé quand $I$ est de longueur $\leq n-1$. (NB Ce théorème est trivial pour $n\leq1$, car alors les catégories $\uHom(\tau,\tau')$ sont toutes ou bien vides, ou bien ponctuelles.)

\msp{page 90}
1) On veut assurer a). Pour ceci, on utilise l'hypothèse de récurrence, pour remplacer un groupe de \texttt{[mot illisible]} flèches critiques consécutives de même signe, par une seule flèche critique. Cette méthode ne foire que si les flèches critiques sont toutes de même signe $\varepsilon$ et incluent $\ul{u}_1$ et $\ul{u}_n$. Mais l'hypothèse de récurrence appliquée à $I=[1,n-1]\subseteq I_\tau$ nous ramène au cas où il n'y a que deux flèches critiques $\ul{u}_1$ et $\ul{u}_n$. Si le premier intervalle critique n'est pas réduit à $\Delta_\varepsilon$, i.e. soit de longueur 2, alors l'hypothèse de récurrence appliquée à $I=[1,n]\subseteq I_\tau$ nous ramène au cas où les deux flèches critiques sont $\ul{u}_1$ et $\ul{u}_2$, et on est
\printedpage{58}
% head: A. Grothendieck : Dérivateurs, Ch. VII, Édition des Universités de Montpellier II et Paris VII
% foot: Version 22 décembre 2001
\msp{page 91}
alors ramené au cas où $\boxed{I=\Delta^2_\varepsilon,\ \tau'=\Delta^2_\varepsilon.}$ Sinon, au cas où $\boxed{\tau=\Delta_\varepsilon\Delta^{l_2}_{-\varepsilon}\Delta_\varepsilon,\ \tau'=\Delta^2_\varepsilon,}$ et où les deux flèches critiques sont encore la première et la dernière, et $\tau'=\Delta^2_\varepsilon$. Ce sont là les deux cas irréductibles, qu'il faudra traiter à part.

Ces cas admis, on a donc moyen pour définir $F$ de façon à satisfaire a). Ensuite on veut :

2) Assurer que la suite des intervalles critiques $I_{j_1}$, $I_{j_2}$, \dots, $I_{j_{r'}}$ qui contient les flèches critiques $\ul{u}_{i_1}$, \dots, $\ul{u}_{i_{r'}}$, soit la suite $I_1$, \dots, $I_{r'}$. C'est encore immédiat par l'hypothèse de récurrence, le seul cas irréductible étant celui où $r'=1$ (une seule flèche critique), et où il s'agit de faire descendre cette flèche critique jusqu'à $\ul{u}_1$. Si $I_1$ est de longueur $>1$, alors on fait descendre en $\ul{u}_2$ par hypothèse de récurrence, et le cas irréductible est celui où $\boxed{\tau=\Delta^2_\varepsilon,\ \tau'=\Delta_\varepsilon,}$ une flèche critique $\ul{u}_2$, la faire descendre en $\ul{u}_1$. Si par contre $I_1$ \texttt{[est]} de longueur 1, le cas irréductible est $\boxed{\tau=\Delta_\varepsilon\Delta^{l_2}_{-\varepsilon}\Delta_\varepsilon,\ \tau'=\Delta_\varepsilon,}$ une flèche critique $\ul{u}_{l_2+2}$, la faire descendre en $\ul{u}_1$.

Sous réserve de traiter les deux nouveaux cas irréductibles, on peut déformer $F$ en un foncteur qui satisfait à a) et à c). Reste à normaliser chacun des $F_\alpha:I_\alpha\longrightarrow I'_\alpha$ ($\alpha\in[1,r']$). Donc on est ramené à $\Delta^n_\varepsilon\longrightarrow\Delta_\varepsilon$ : une seule flèche critique $\ul{u}_\alpha$, $\alpha\in[1,n]$, à faire descendre en $\ul{u}_1$. Si $n\geq2$, par hypothèse de récurrence elle descend en $\ul{u}_2$, cas supposé traité en $2^o$). Donc
\msp{page 92}
on gagne, sous réserve de traiter les 4 cas irréductibles.

A) Cas $\tau=\Delta^2_\varepsilon$, $\tau'=\Delta^2_\varepsilon$ ou $\Delta_\varepsilon$.

La catégorie $\uHom(\Delta^2_\varepsilon,\Delta_\varepsilon)$ a deux éléments, les morphismes $\Delta^2_\varepsilon\longrightarrow\Delta_\varepsilon$ correspondent aux flèches de transition de $\Delta^2_\varepsilon$ ($=\mathbf{\Delta}^+_2$ ou $\mathbf{\Delta}^-_2$). Dans le cas $\varepsilon=+1$, ce sont les deux applications de $\{0,1,2\}$ dans $\{0,1\}$, qui appliquent 0 en 0, 2 en 1, et 1 soit en 1 ($F_0$), soit en 1 \texttt{[plutôt en 0]} ($F_1$).
\[
\begin{tikzcd}
0 \arrow[r] & 1 \arrow[r] & 2
\end{tikzcd}
\]
\[
\begin{tikzcd}
0 \arrow[r] & 1
\end{tikzcd}
\]

On a évidemment une flèche $F_1\longrightarrow F_0$ (le voyageur pressé arrive plus vite \ldots). Similaire pour $\varepsilon=-1$,
\[
\uHom(\Delta^2_\varepsilon,\Delta_\varepsilon)\simeq\mathbf{\Delta}^1,\qquad\text{catégorie 0-connexe.}
\]

La catégorie $\uHom(\Delta^2_\varepsilon,\Delta^2_\varepsilon)$ a trois éléments, s'identifiant à l'ensemble des parties non vide de l'ensemble $\mathrm{Tr}(\tau)$ à deux éléments des flèches de transition $\ul{u}_1$, $\ul{u}_2$ de $\tau$ : il y a $F_{\ul{u}_1}=F_0$, $F_{\ul{u}_2}$, $F_{\ul{u}_1,\ul{u}_2}$. Si $\varepsilon=+1$, il s'agit des applications croissantes (sans plus) de $\{0,1,2\}$ dans $\{0,1,2\}$, appliquant 0 en 0, 2 en 2, et 1 en ce qu'elle voudra. % [sic] parties non vide -> parties non vides, agreement

% ---- from Der-ChVII-456_p061-065.tex
\printedpage{59}
% head: A. Grothendieck : Dérivateurs, Ch. VII, Édition des Universités de Montpellier II et Paris VII
% foot: Version 22 décembre 2001

\begin{tabular}{lll}
Donc & $1\longmapsto 2$ & correspond à \quad $F_0$\\[1ex]
 & $1\longmapsto 1$ & correspond à \quad $F_1$\\[1ex]
 & $1\longmapsto 0$ & correspond à \quad $F_2$
\end{tabular}

et on a bien sur $F_2\longrightarrow F_1\longrightarrow F_0$, % [sic] bien sur -> bien sûr, accent missing
\[
\uHom(\Delta^2_\varepsilon,\Delta^2_\varepsilon)\simeq\boldsymbol{\Delta}^2,\qquad\quad\text{catégorie 0-connexe.}
\]

Dans ces deux cas, deux éléments de $\mathrm{Hom}(\tau,\tau')$ sont reliés par une homotopie \textit{simple.}

\msp{page 93}

B) Cas $\tau=\Delta_\varepsilon\Delta^l_{-\varepsilon}\Delta_\varepsilon$, $\tau'=\Delta^2_\varepsilon$ ou $\Delta_\varepsilon$.

Notons que l'on a, pour $\tau,\tau'\in\uTy$ quelconque, un isomorphisme de catégories
\[
\uHom(\tau,\tau')\longrightarrow\uHom(\tau'^o,\tau^o),
\]
exprimant simplement que
\[
\tau\longmapsto\tau^o
\]
est un \texttt{[2]}-foncteur de $\uTy$ dans $\uTy$, et même un automorphisme de cette \texttt{[2]}-catégorie. Donc dire que $\uHom(\tau,\tau')$ est connexe revient à dire que $\uHom(\tau^o,\tau'^o)$ l'est. Donc on peut supposer ci-dessus que $\varepsilon=+1$ (et pareil dans A!).
\[
\begin{aligned}
\tau :\quad & \longrightarrow\ \underbrace{\longleftarrow\ \longleftarrow\ \cdots\ \longleftarrow}_{l}\ \longrightarrow\\
\tau' :\quad & \longrightarrow\ \longrightarrow\quad\text{ou}\quad\longrightarrow\ .
\end{aligned}
\]

Mais tous les morphismes $\tau\longrightarrow\tau'$ sont triviaux sur l'intervalle critique $I_2$ (rétrograde), ce qui implique que
\[
\uHom(\tau,\tau')\xleftarrow{\ \sim\ }\uHom(\Delta^2_\varepsilon,\tau'),
\]
via le foncteur de contraction canonique
\[
I_\tau\longrightarrow I_{\Delta^2_\varepsilon}=\boldsymbol{\Delta}^2.
\]

Or dans A on a vu que $\uHom(\Delta^2_\varepsilon,\tau')$ est connexe (pour $\tau'$ de la forme $\Delta_\varepsilon$ ou $\Delta^2_\varepsilon$).

On a gagné!

\newpage
\printedpage{60}
% head: A. Grothendieck : Dérivateurs, Ch. VII, Édition des Universités de Montpellier II et Paris VII
% foot: Version 22 décembre 2001

\msp{page 94}

\section*{9\quad Théorème d'homotopie pour les catégories de chemins}

Soient $X$ une catégorie, $\Sigma\subseteq\mathrm{Fl}(X)$ \texttt{[un]} ensemble admissible de flèches dans $X$, $a,b\in\mathrm{Ob}\,X$. Considérons pour $\tau\in\mathrm{Ty}$ variable, les catégories
\[
\uCh_{\Sigma,\tau}(X),\qquad\quad\uCh_{\Sigma,\tau}(X;a,b),
\]
qui forment les fibres des catégories Cat-fibrées
\[
\begin{tikzcd}[column sep=small]
\uCh_\Sigma(X) \arrow[dr] & & \uCh_\Sigma(X;a,b) \arrow[ll, "\text{incl.}"'] \arrow[dl]\\
 & \uTy &
\end{tikzcd}
\]
sur la catégorie $\uTy$. Ainsi, pour $\tau,\tau'\in\mathrm{Ty}$, et pour toute
\[
\varphi:\tau\longrightarrow\tau',
\]
on a un carré commutatif d'homomorphismes
\[
\begin{tikzcd}
\uCh_{\Sigma,\tau'}(X) \arrow[d, "\varphi^*_X"'] & \uCh_{\Sigma,\tau'}(X;a,b) \arrow[l, "\text{incl.}"'] \arrow[d, "\varphi^*_{a,b}"]\\
\uCh_{\Sigma,\tau}(X) & \uCh_{\Sigma,\tau}(X;a,b), \arrow[l, "\text{incl.}"']
\end{tikzcd}
\]
avec transitivité pour $\varphi$ variable (composition $\psi\circ\varphi$ de flèches dans $\uTy$).

En fait, pour $\varphi$ variable dans $\uHom(\tau,\tau')$, les foncteurs images inverses $\varphi^*_X$, $\varphi^*_{a,b}$ dépendent \textit{fonctoriellement} de $\varphi$; i.e. on a des \textit{foncteurs}
\[
\begin{tikzcd}[row sep=small]
 & \uHom(\uCh_{\Sigma,\tau'}(X),\uCh_{\Sigma,\tau}(X))\\
\uHom(\tau,\tau') \arrow[ur] \arrow[dr] & \\
 & \uHom(\uCh_{\Sigma,\tau'}(X;a,b),\uCh_{\Sigma,\tau}(X;a,b))
\end{tikzcd}
\tag*{$(*)$}
\]

\newpage
\printedpage{61}
% head: A. Grothendieck : Dérivateurs, Ch. VII, Édition des Universités de Montpellier II et Paris VII
% foot: Version 22 décembre 2001

\msp{page 95}

Cela provient de l'inclusion
\[
\uHom(\tau,\tau')\hookrightarrow\uHom^{\varepsilon}(I_\tau,I_{\tau'})
\]
et de la description des catégories $\uCh_\tau$, $\uCh_{\tau'}$ comme catégories de foncteurs, telle
\[
\uCh_{\Sigma,\tau}(X)\simeq\uHom_{\Sigma_\tau,\Sigma}(I_\tau,X),
\]
où $\Sigma_\tau$ désigne l'ensemble des flèches dans $I_\tau$ qui sont soit inversibles, soit \textit{rétrogrades.} (Cet ensemble de flèches est admissible, car stable par composition.) On désigne par $\uHom_{\Sigma_\tau,\Sigma}$ la sous-catégorie pleine de $\uHom$ formée des foncteurs qui appliquent $\Sigma_\tau$ dans $\Sigma$. Or de façon générale, si
\[
(I,\Sigma_I),\ \ (I',\Sigma_{I'}),\ \ (X,\Sigma)
\]
sont trois catégories, munies d'ensembles admissibles de flèches, considérant les foncteurs $\varphi:I\longrightarrow I'$ avec $\varphi(\Sigma_I)\subseteq\Sigma_{I'}$, alors le foncteur familier
\[
\uHom(I,I')\longrightarrow\uHom(\uHom(I',X),\uHom(I,X))
\]
(transposé du foncteur de composition
\[
\begin{matrix}
\uHom(I,I') & \times & \uHom(I',X) & \longrightarrow & \uHom(I,X) & \\
(\varphi & , & F') & \longmapsto & F'\circ\varphi & ),
\end{matrix}
\]
se généralise en un foncteur
\[
\uHom_{\Sigma_I,\Sigma_{I'}}(I,I')\longrightarrow\uHom(\uHom_{\Sigma_{I'},\Sigma}(I',X),\uHom_{\Sigma_I,\Sigma}(I,X)),
\tag*{$(*)$}
\]
transposé du foncteur de composition
\[
\uHom_{\Sigma_I,\Sigma_{I'}}(I,I')\times\uHom_{\Sigma_{I'},\Sigma}(I',X)\longrightarrow\uHom_{\Sigma_I,\Sigma}(I,X)
\]
induit par $(*)$.

\msp{page 96}

Comme toute
\[
\varphi:\tau\longrightarrow\tau'
\]
induit
\[
\varphi:I_\tau\longrightarrow I_{\tau'},
\]
transformant flèches rétrogrades en flèches identiques ou rétrogrades, i.e.
\[
\varphi(\Sigma_\tau)\subseteq\Sigma_{\tau'},
\]
\printedpage{62}
% head: A. Grothendieck : Dérivateurs, Ch. VII, Édition des Universités de Montpellier II et Paris VII
% foot: Version 22 décembre 2001
on trouve le premier des deux foncteurs $(*)$ annoncés page 94, qui est associé par transposition au foncteur $\gamma$
\[
\begin{tikzcd}[column sep=large]
\uHom(\tau,\tau')\times\ \underbrace{\uCh_{\Sigma,\tau'}(X)}_{\simeq\,\uHom_{\Sigma_{\tau'},\Sigma}(I_{\tau'},X)} \arrow[r, "\gamma"] \arrow[d, "\text{can.}"'] & \uCh_{\Sigma,\tau}(X) \arrow[d, equal]\\
\uHom_{\Sigma_\tau,\Sigma_{\tau'}}(I_\tau,I_{\tau'})\times\uHom_{\Sigma_{\tau'},\Sigma}(I_{\tau'},X) \arrow[r, dotted, "\gamma_0"] & \uHom_{\Sigma_\tau,\Sigma}(I_\tau,X)
\end{tikzcd}
\tag*{$(*)$}
\]
déduit du foncteur composition $\gamma_0$ dans ce diagramme : $\gamma=\gamma_0\circ\mathrm{can}.$ D'autre part, ce foncteur de composition $\gamma$ induit visiblement un foncteur sur les sous-catégories de chemins à source et but fixés
\[
\uHom(\tau,\tau')\times\uCh_{\Sigma,\tau'}(X;a,b)\longrightarrow\uCh_{\Sigma,\tau}(X;a,b),
\]
d'où la deuxième des deux flèches annoncées dans $(*)$ p. 94.

Cela montre donc que si
\[
\varphi_1,\varphi_2:\tau\rightrightarrows\tau'
\]
sont deux morphismes de $\tau$ dans $\tau'$ homotopes dans $\uHom(\tau,\tau')$, alors les foncteurs induits entre catégories de $\Sigma$-chemins ou entre catégories de $\Sigma$-chemins à extrémités précisés : % [sic] précisés -> précisées, agreement with extrémités

\msp{page 97}
\[
\begin{aligned}
\varphi^*_{1\,X},\varphi^*_{2\,X} : && \uCh_{\Sigma,\tau'}(X) &\quad\rightrightarrows\quad\uCh_{\Sigma,\tau}(X),\\
(\varphi_1)^*_{a,b},(\varphi_2)^*_{a,b} : && \uCh_{\Sigma,\tau'}(X;a,b) &\quad\rightrightarrows\quad\uCh_{\Sigma,\tau}(X;a,b)
\end{aligned}
\]
sont également homotopes.

Or le lemme d'homotopie nous assure que l'hypothèse $\varphi_1\sim\varphi_2$ est \textit{toujours} satisfaite. Donc :

\textbf{Théorème d'homotopie sur les opérations universelles sur les chemins.} \textit{Soient $X$, $\Sigma$, $a$, $b$ comme dessus, et $\tau,\tau'\in\mathrm{Ty}$ telles que $\mathrm{Hom}(\tau,\tau')\neq\emptyset$. Alors les foncteurs $\varphi^*_X$, $\varphi^*_{X;a,b}$ induits par un $\varphi\in\mathrm{Hom}(\tau,\tau')$}
\[
\left\{
\begin{aligned}
\varphi^*_X : && \uCh_{\Sigma,\tau'}(X) &\quad\longrightarrow\quad\uCh_{\Sigma,\tau}(X)\\
\varphi^*_{X;a,b} : && \uCh_{\Sigma,\tau'}(X;a,b) &\quad\longrightarrow\quad\uCh_{\Sigma,\tau}(X;a,b)
\end{aligned}
\right.
\]
\textit{ne dépendent pas de $\varphi$,} à homotopie près.

On peut préciser cet énoncé, de telle façon que l'énoncé pour les $\varphi^*_{X;a,b}$ apparaisse, comme il devrait, comme corollaire pour celui sur $\varphi^*_X$. Pour ceci, on va regarder les catégories de chemins dans $X$ comme catégories \textit{au dessus de $X$,} par les foncteurs `source et but' :
\[
\begin{tikzcd}[column sep=small]
\uCh_{\Sigma,\tau}(X) \arrow[dr, "(\sigma_\tau{,}\beta_\tau)"'] & & \uCh_{\Sigma,\tau'}(X) \arrow[dl, "(\sigma_{\tau'}{,}\beta_{\tau'})"]\\
 & X\times X &
\end{tikzcd}
\]

\newpage
\printedpage{63}
% head: A. Grothendieck : Dérivateurs, Ch. VII, Édition des Universités de Montpellier II et Paris VII
% foot: Version 22 décembre 2001

\msp{page 98}

Alors le foncteur $(*)$ (p. 94) se factorise par la sous-catégorie $\uHom_{\mathrm{Cat}/X\times X}(\cdot,\cdot)$ de $\uHom_{\mathrm{Cat}}(\cdot,\cdot)$ formée des foncteurs qui commutent aux foncteurs de projection dans $X\times X$, et les homomorphismes entre tels foncteurs qui s'effectuent par des flèches `verticales' de la catégorie but (i.e. par des flèches au dessus de flèches identiques de la catégorie base $X\times X$) :
\[
\uHom(\tau,\tau')\longrightarrow\uHom_{\mathrm{Cat}/X\times X}(\uCh_{\Sigma,\tau'}(X),\uCh_{\Sigma,\tau}(X)).
\]

Donc dans l'énoncé du théorème d'homotopie, on peut préciser que l'homotopie se fait dans cette sous-catégorie, i.e. s'effectue par une suite d'homotopies simples qui sont des morphismes de foncteurs entre catégories au dessus de $X\times X$. Cela implique alors l'énoncé d'homotopie sur les $\varphi^*_{X;a,b}$.

% ---- from Der-ChVII-456_p066-070.tex
\leavevmode\printedpage{64}
% head: A. Grothendieck : Dérivateurs, Ch. VII, Édition des Universités de Montpellier II et Paris VII
\msp{page 99}

\section*{10\quad Question de comparaison pour les catégories \textmd{Cat}-fibrées satisfaisant au ‘théorème d'homotopie' (base filtrante)}

Nous allons généraliser maintenant au cas d'une catégorie (Cat)-fibrée sur une autre
\[
\begin{tikzcd}
C \arrow[d, "\pi"'] \\
T,
\end{tikzcd}
\]
en pensant surtout au cas $T = \uTy$, $C = \uCh_{\Sigma}(X;a,b)$. L'hypothèse essentielle est la suivante :

\noindent\begin{tabular}{@{}p{0.3\textwidth}p{0.64\textwidth}@{}}
(Hypothèse d'homotopie) & Si $t,t' \in \Ob T$ sont tels que $\Hom(t,t') \neq \emptyset$, alors les foncteurs $\varphi^{*} : C_{t'} \longrightarrow C_{t}$ associés aux $\varphi \in \Hom(t,t')$ sont homotopes.
\end{tabular}

Ils définissent donc une seule et même flèche dans Cat\,hot. Je désigne par $T_{\mathrm{ord}}$ ou $\To$ l'ensemble $\Ob(T)$, avec comme relation d'ordre celle déduite de la structure de catégorie :
\[
t \leq t' \qquad\Longleftrightarrow\qquad \Hom_{T}(t,t') \neq \emptyset,
\]
de sorte qu'on a un foncteur
\[
T \longrightarrow \To
\]
qui est \textit{universel} pour les foncteurs de $T$ dans des catégories ‘ordonnées' (i.e. où toute flèche est monomorphique \texttt{[non, cf. }$\bullet\rightrightarrows\bullet$\texttt{; plutôt, où toujours }$\#\Hom(a,b) \leq 1$ \texttt{(admettant des isomorphismes non identiques)]}) ou pour les foncteurs $F$ sur $T$ tels que $F$ soit constant sur tout $\Hom(t,t')$. L'hypothèse d'homotopie peut s'exprimer en disant que le foncteur composé
\[
\begin{tikzcd}
T^{o} \arrow[r] & \Cat \arrow[r] & \Cathot \\
t \arrow[r, mapsto] & C_{t} &
\end{tikzcd}
\]
\msp{page 100}
se factorise (évidemment de façon unique) en un foncteur
\[
H_{C} : \To^{o} \longrightarrow \Cathot,
\]
i.e. qu'on a un diagramme commutatif
\printedpage{65}
% head: A. Grothendieck : Dérivateurs, Ch. VII, Édition des Universités de Montpellier II et Paris VII
\[
\begin{tikzcd}
T^{o} \arrow[r, "t\,\longmapsto\, C_{t}"] \arrow[d] & \Cat \arrow[d] \\
\To^{o} \arrow[r, "H_{C}"] & \Cathot.
\end{tikzcd}
\]

Nous allons supposer de plus :

\noindent\begin{tabular}{@{}p{0.36\textwidth}p{0.58\textwidth}@{}}
(Hypothèse de filtration) & L'ensemble ordonné $\To$ associé à $T$ est \textit{filtrant} décroissant (filtrant à gauche), i.e. si $t,t' \in \Ob T$, il existe $t'' \in \Ob T$ et des flèches $t'' \longrightarrow t$, $t'' \longrightarrow t'$.
\end{tabular}

Cette hypothèse est bien satisfaite dans $\uTy$, car pour $\tau,\tau' \in \Ob \uTy = \mathrm{Ty}$, il suffit de prendre $\tau\tau'$, qui s'envoie dans $\tau$ par $\mathrm{id}_{\tau} * p_{\tau'}$, dans $\tau'$ par $p_{\tau} * \mathrm{id}_{\tau'}$, où $p_{\tau} : \tau \longrightarrow 1$ est le morphisme de $\tau$ dans l'objet final.

Ainsi, $H_{C}$ définit un ind-objet de Cat\,hot, que nous allons désigner par $\underline{\underline{C}}$. Si on suppose de plus :

\noindent\begin{tabular}{@{}p{0.36\textwidth}p{0.58\textwidth}@{}}
(Hypothèse de dénombrabilité, toujours satisfaite si $\mathrm{card}\,\To \leq \aleph_{0}$) & $\To$ admet un sous-ensemble \textit{dénombrable} cofinal à gauche.
\end{tabular}

alors cet ind-objet de Cat\,hot peut être défini, de beaucoup de façons, en termes
\msp{page 101}
d'un système inductif dans Cat (pas seulement dans Cat\,hot)
\[
(C_{n})_{n\in\mathbf{N}},
\]
d'où on récupère l'ind-objet $\underline{\underline{C}}$ dans Cat\,hot en passant au système inductif correspondant dans Cat\,hot, puis à l'ind-objet associé.

Pour construire le système inductif $C_{n}$, il suffit de prendre une chaîne infinie
\[
t_{0} \longleftarrow t_{1} \longleftarrow t_{2} \longleftarrow \;\cdots\; \longleftarrow t_{n} \longleftarrow t_{n+1} \longleftarrow \;\cdots
\]
de morphismes dans $T$, telle que l'ensemble des $t_{i}$ soit cofinal à gauche dans $\To$. On va supposer de plus $C$ \textit{strictement fibrée} sur $T$ (i.e. provenant d'un véritable foncteur $T^{o} \longrightarrow \Cat$, i.e. on peut choisir les foncteurs images inverses
\[
\varphi^{*} : C_{t'} \longrightarrow C_{t}
\]
de façon transitive stricte, pas seulement transitive à isomorphisme près) - cas auquel on peut toujours se ramener d'ailleurs par l'astuce de \textsc{Giraud}. Cette hypothèse est d'ailleurs vérifiée dans le cas $\uCh_{\Sigma}(X;a,b)$ sur $\uTy$. Ainsi, les
\[
(C_{t_{i}})_{i\in\mathbf{N}}
\]
forment un système inductif de catégories indexées par $\mathbf{N}$, le morphisme de transition
\[
C_{n} \longrightarrow C_{n'} \qquad \text{i.e. } C_{t_{n}} \longrightarrow C_{t_{n'}}
\]
\printedpage{66}
% head: A. Grothendieck : Dérivateurs, Ch. VII, Édition des Universités de Montpellier II et Paris VII
\msp{page 102}
étant le foncteur image inverse $\varphi^{*}_{n,n'}$, où
\[
\varphi_{n,n'} : t_{n'} \longrightarrow t_{n}
\]
est le composé des $t_{n'} \longrightarrow t_{n'-1} \longrightarrow \;\cdots\; \longrightarrow t_{n}$. (Si on n'avait pas transitivité stricte, on définirait $\varphi^{*}_{n,n'}$ comme le composé des foncteurs image inverse par des flèches de transition successives $t_{n'} \longrightarrow t_{n'-1} \longrightarrow \;\cdots\; \longrightarrow t_{n+1} \longrightarrow t_{n}$, ce ne serait $\varphi^{*}_{n,n'}$ qu'à isomorphisme canonique près.)

Soit
\[
C_{\infty} = \varinjlim_{n} C_{n} = \varinjlim_{n} C_{t_{n}}
\]
la catégorie $\varinjlim$ . (Elle sera surtout sympathique si on a pu choisir le système $(t_{n})_{n\in\mathbf{N}}$ de telle façon que les foncteurs de transition
\[
C_{t_{n}} \longrightarrow C_{t_{n-1}}
\]
soient pleinement fidèles, et peut-être même \texttt{[des]} monomorphismes, de telle façon que les $C_{n}$ s'identifient à des sous-catégories pleines de leur limite inductive $C_{\infty}$.) Je voudrais comprendre les relations précises entre les trois genres d'objets suivants :

\begin{itemize}[leftmargin=2.5em]
\item[\circled{a}] $C$ elle-même (comme objet dans Cat, pas de $\Cat/T$), et surtout les objets de Cat\,hot et de Hot qui lui correspondent.

\item[\circled{b}] \msp{page 103}$C_{\infty}$, et les objets de Cat\,hot et de Hot qui lui correspondent. Dans quelle mesure ces deniers % [sic] deniers -> derniers, missing letter
sont-ils indépendant % [sic] indépendant -> indépendants, agreement with the plural subject
du choix du système projectif cofinal à gauche $(t_{n})_{n\in\mathbf{N}}$ dans $T$? Quel est le rôle de $C_{\infty}$? Est-elle canoniquement isomorphe à $C$ dans Cat\,hot, ou dans Hot? (Auquel cas la question de son ‘rôle' me semblerait résolue!)

\item[\circled{c}] Le ind-objet % [sic] Le ind-objet -> L'ind-objet, missing elision
$H_{C}$ dans Cat\,hot, et l'ind-objet de Hot qu'il définit.

Je note que l'ind-objet $H_{C}$, et aussi les catégories $C_{\infty}$ construites comme dessus, ne changent pas quand on remplace $T$ par une sous-catégorie $T'$ cofinale à gauche, ou aussi par une catégorie de la forme $T/t_{0}$ ($t_{0} \in \Ob T$), qui donne lieu à l'ensemble ordonnée % [sic] ordonnée -> ordonné, agreement with "ensemble"
associé
\[
(T/t)_{\mathrm{o}} = \To/t = \{t \in \To \mid t \leq t_{0}\}
\]
\texttt{[plutôt} $(T/t_{0})_{\mathrm{o}} = \To/t_{0} = \{t \in \To \mid t \leq t_{0}\}$\texttt{]}. Si la relation entre les objets dans a), et ceux dans b) et dans c), est très étroite, on en concluerait % [sic] concluerait -> conclurait, spelling of the conditional
que le ‘type' de $C$ dans Cat\,hot, ou dans Hot suivant le cas, ne change pas en remplaçant $C$ par
\printedpage{67}
% head: A. Grothendieck : Dérivateurs, Ch. VII, Édition des Universités de Montpellier II et Paris VII
\msp{page 104}
$C \times_{T} T'$ ou par $C/t_{0} \simeq C \times_{T} T/t_{0}$. Bien mieux, elle ne changerait pas en remplaçant $T$ par $(\mathbf{N})^{o}$, la catégorie (ordonnée) opposée à $\mathbf{N}$ (ayant donc un unique objet final), et la catégorie fibrée $C$ sur $T$ par la catégorie $C_{\underline{t}}$ induite par $C$ sur $(\mathbf{N})^{o}$, via le foncteur
\[
\underline{t} : (\mathbf{N})^{o} \longrightarrow T
\]
défini par le système projectif $\underline{t} = (t_{n})_{n\in\mathbf{N}}$ dans $T$. Cela signifierait donc que pour comprendre le type homotopique (au sens strict, i.e. dans Cat\,hot, ou large, i.e. dans Hot) de $C$, on peut le ‘nettoyer' de façon draconienne, la base $T$ étant vue comme une sorte de gangue un peu vaseuse, ou du moins pléthorique.

\item[\circled{d}] De toutes façons, si $T'$ est une autre catégorie filtrante à gauche, et
\[
\alpha : T' \longrightarrow T
\]
un foncteur cofinal à gauche, de sorte que $C$ sur $T$, et $C' = C \times_{T} T'$ sur $T'$, définissent essentiellement le ‘même' objet de Ind\,Cat\,hot, et que la catégorie $C'_{\infty}$ pour $C'$ sur $T'$ fait partie des catégories $C_{\infty}$ pour $C$ sur $T$ - on peut se demander
\msp{page 105}
si le foncteur canonique
\[
C' \;\longrightarrow\; C
\]
est également toujours un isomorphisme dans Cat\,hot, ou sinon, du moins dans Hot.
\end{itemize}

Un intérêt particulier s'attache à $\pi_{0}(C)$, et de plus à sa variance pour $C$ sur $T$ variable de la façon suivante. Si $S$ est dans Cat, on regarde les catégories
\[
\uHom_{\Cat}(S,C_{t}) \qquad (t \in T),
\]
on trouve un foncteur
\[
t \longmapsto \uHom_{\Cat}(S,C_{t}), \qquad T^{o} \longrightarrow (\Cat),
\]
d'où une catégorie fibrée sur $T$, que je vais désigner par
\[
\sHom_{T}(S_{T},C)
\]
(Sauf erreur, elle s'interprète comme la catégorie fibrée sur $T$ correspondante au champ sur $\mathrm{Top}(T) = T^{\wedge}$ du morphisme de $T$-champs de $S_{T}$ dans $C$, où $S_{T} = S \times T$ (considérée comme catégorie fibrée sur $T$ à fibre essentiellement constante $S$) dans $C$. Ici j'identifie catégorie fibrée sur $T$, et champs sur le topos $T^{\wedge}$.) Il est immédiat que la propriété d'homotopie (p. 99) est stable par passage de $C$ à $\sHom_{T}(S_{T},C)$. On s'intéresse à la variation de
\[
\pi_{0}(\sHom_{T}(S_{T},C))
\]
\printedpage{68}
% head: A. Grothendieck : Dérivateurs, Ch. VII, Édition des Universités de Montpellier II et Paris VII
\msp{page 106}
comme contrafoncteur en $S$,
\[
H : S \longmapsto \pi_{0}\sHom_{T}(S_{T},C) \qquad (\Cat)^{o} \longrightarrow (\mathrm{Ens}),
\]
et à une éventuelle expression de ce contrafoncteur $H$ en termes de l'ind-objet $H_{C}$ de Cat\,hot, par la formule hypothétique
\[
H(S) \simeq \Hom_{\mathrm{Ind}(\Cathot)}(\gamma(S),H_{C}),
\]
où
\[
\gamma : \Cat \longrightarrow \Cathot
\]
est le foncteur canonique. Si on avait un tel isomorphisme canonique fonctoriel en $S$, cela donnerait une description (à isomorphisme unique près) de l'ind-objet $H_{C}$ en termes du foncteur précédent $H$ sur Cat.

Il me vient un doute ici, en vue de l'exemple $T = \uTy$, $C = \uCh_{\Sigma}(X;a,b)$. Les trois types d'objets a) b) c) ont le même $\pi_{0}$, et il est sûr que ‘c'est le bon'. D'autre part, pour une catégorie paramètre variable $S$, il est sûr aussi que le bon foncteur en $S$ est
\[
S \longmapsto \pi_{0}(\uHom_{T}(S_{T},C)) \simeq \Hom_{\mathrm{Ind}\,\Cathot}(S,H_{C}),
\]
isomorphe aussi, \textit{si $S$ est finie,} à
\[
\Hom_{\Cathot}(S,C_{\infty}),
\]
ce qui montre que, du moins pour des arguments $S \in (\Cat)$ qui sont des catégories \textit{finies}, les objets
\msp{page 107}
$C_{\infty}$ et $H_{C}$ sont ‘bons', il nous permettent % [sic] il nous permettent -> ils nous permettent, agreement with the plural subject
de décrire un foncteur important qui domine la situation envisagée. Moyennant une hypothèse très forte, mais satisfaite dans les cas importants, on peut dans ces deux formules remplacer Cat\,hot par Hot.

Par contre, sauf pour le cas où $S$ est réduit à un point et où $\Hom_{\Cathot}$ ou $\Hom_{\Hot}$ de $S$ dans $C$ n'est autre que $\pi_{0}(C)$ ($\simeq \pi_{0}(C_{\infty}) \simeq \pi_{0}(H_{C})$), donc effectivement ‘bon', je ne sais pas si
\[
\underbrace{\Hom_{\Cathot}(S,C)}_{=\,\pi_{0}\uHom(S,C)} \qquad \text{ou} \qquad \Hom_{\Hot}(S,C)
\]
est lié de façon raisonnable aux ensembles précédents (déterminés par $H_{C}$, resp. par $C_{\infty}$), p. ex. s'ils lui sont isomorphes de façon naturelle, et s'ils ont une signification géométrico-topologique raisonnable. Je commence à suspecter que \textit{non} - donc qu'on aurait ici un cas où une relation de type très général concernant un $\pi_{0}$, ici donc
\[
\pi_{0}(C) \simeq \pi_{0}(C_{\infty}) \simeq \pi_{0}(H_{C}),
\]
ne se prolonge pas en un isomorphisme sur les $\pi_{i}$ supérieurs, et en une propriété de Hot-équivalence.

% ---- from Der-ChVII-456_p071-075.tex
\printedpage{69}
% head: A. Grothendieck : Dérivateurs, Ch. VII, Édition des Universités de Montpellier II et Paris VII

\msp{page 108}
Je ne trouve d'ailleurs aucune flèche directe entre $C$ et l'un des objets $C_{\infty}$, $H_{C}$; tout ce que je vois dans ce sens, c'est
\[
\begin{tikzcd}[column sep=large, row sep=huge]
\CNo \arrow[r, "{\substack{\text{foncteur de}\\ \text{localisation}}}"] \arrow[d] & C_{\infty} \\
C, &
\end{tikzcd}
\]
où la flèche verticale est la flèche de projection de $\CNo = C \times_{T} \mathbf{N}^{o}$ (avec $\mathbf{N}^{o} \xrightarrow{\ t\ } T$). Ceci n'établit une relation topologique étroite entre $C$ (en tant qu'objet de Cat\,hot ou de Hot simplement) et $C_{\infty}$, que si les deux flèches précédentes sont tout au moins des Hot-équivalences. Et le cas-type (problématique) où la question de la relation entre $C$ et $C_{\infty}$ se pose, est donc celui où $T = \mathbf{N}^{o}$, i.e. celui d'un système inductif quelconque de (petites) catégories, indexé par $\mathbf{N}$, et interprété comme défini par une catégorie fibrée $C = \CNo$ sur $\mathbf{N}^{o}$. Dans ce cas particulier, on a du moins une flèche canonique
\[
C \longrightarrow C_{\infty}
\]
sur laquelle on peut se poser la question si c'est
\msp{page 109}
une Hot-équivalence. Sinon, cela prouvera que dans le cas général, le Hot-type de $C$ (et a fortiori sa classe d'isomorphie dans Cat\,hot) n'a pas de sens géométrique intéressant. C'est à dire que $C$ est intéressante uniquement \textit{en tant que catégorie fibrée sur $T$}, permettant notamment de construire l'ind-objet $H_{T}$ de Cat\,hot ou de Hot, mais pas en tant que `catégorie nue'.

Mais je viens de me convaincre que $\CNo \longrightarrow C_{\infty}$ est bien une Hot-équivalence. Donc la flèche $\CNo \longrightarrow C$ définit une flèche \textit{dans} Hot
\[
C_{\infty} \longrightarrow C
\]
(en composant avec l'inverse dans Hot de la précédente). Est-ce un isomorphisme? Donc pour cette question, on est ramené à la question pour $\CNo \longrightarrow C$. Mais si on prend $C = T$, cela pose déjà la même question pour $\mathbf{N}^{o} \longrightarrow T$, donc la question si $T$ est asphérique, sous la seule hypothèse que $T_{\mathrm{o}}$ soit filtrante à gauche et dénombrable à l'infini (à gauche). Mais bien sûr c'est faux (p. ex. si $T$ n'a qu'un seul objet, il n'est pas pour autant asphérique !).

Cela me convainc donc qu'à défaut
% foot: Version 22 décembre 2001
\printedpage{70}
% head: A. Grothendieck : Dérivateurs, Ch. VII, Édition des Universités de Montpellier II et Paris VII
\msp{page 110}
d'hypothèse supplémentaire draconienne sur $T$, la catégorie $C$ `nue' n'offre pas d'intérêt particulier, à part son $\pi_{0}$ (qui est donné aussi par $C_{\infty}$, et par $H_{C}$). Il n'est pas dit, certes, que \underline{Ty} ne satisfasse pas à ces `conditions draconiennes'. Mais comme $C$ `nue' ne s'insère pas dans un formalisme sympathique (p.ex. on n'a pas d'interprétation géométrique de $\mathrm{Hom}_{\mathrm{Cat\,hot}}(S,C)$, ou de $\mathrm{Hom}_{\mathrm{Hot}}(S,C)$), je crois que par la suite je peux en faire abstraction. Exit a) !

\begin{flushright}
{\footnotesize 5.11.90}
\end{flushright}

J'ai quand même passé la journée entière hier à essayer de comprendre la relation de $C$ et de son type dans Hot, à celui de $C_{\infty}$. En effet, comme on a finalement une flèche dans Hot
\[
C_{\infty} \longrightarrow C,
\]
déduite du diagramme dans (Hot)
\[
\begin{tikzcd}[column sep=huge, row sep=huge]
\CNo \arrow[r, "\sim"] \arrow[d] & C_{\infty} \\
C, &
\end{tikzcd}
\]
de prouver que $C_{\infty} \longrightarrow C$ est un isomorphisme dans Hot
\msp{page 111}
semble la façon la plus simple, et peut-être unique, de prouver que l'image de $C_{\infty}$ dans Hot ne dépend pas, à isomorphisme unique près, des choix faits pour construire $C_{\infty}$. Sans compter que dans le cas motivant, $C = \Ch_{\Sigma}(X)$ ou $\Ch_{\Sigma}(X;a,b)$, $T = \Ty$, c'est la catégorie des chemins $C$ qui donne lieu à une loi de composition des chemins impeccable, fonctorielle, associative et tout, alors qu'il n'en est rien pour $C_{\infty}$. On ne trouve pas une loi de composition
\[
C_{\infty} \times C_{\infty} \longrightarrow C_{\infty}
\]
dans Cat, ni même dans Cat\,hot ou dans Hot - ou plutôt, si on arrive pourtant à construire une flèche $C_{\infty} \times C_{\infty} \longrightarrow C_{\infty}$ dans Cat, c'est encore laborieusement à coups de choix $\pm$ arbitraires, cran par cran, et sans savoir si ce qu'on obtient a une propriété d'unicité, ne serait ce que dans Hot; on ne trouve une interprétation géométrique tangible de cette composition que dans $\Hotz$ (où on a une véritable loi de groupoïde), ou dans $H_{T}$ en tant qu'objet de Ind\,Cat\,hot
% foot: Version 22 décembre 2001
\printedpage{71}
% head: A. Grothendieck : Dérivateurs, Ch. VII, Édition des Universités de Montpellier II et Paris VII
\msp{page 112}
ou de Ind\,Hot. (Donc la limite inductive dans $\Hotz$ n'est autre que l'élément de $\Hotz$ défini par $C_{\infty}$.)

Mes efforts pour prouver que
\[
C_{\infty} \longrightarrow C
\]
était un isomorphisme dans Hot, ou ne serait-ce qu'un 1-isomorphisme (i.e. un isomorphisme non seulement pour les $\pi_{0}$, ce qui est toujours vrai, mais aussi pour les $\pi_{1}$ de toutes les composantes), ont échoués. J'ai trouvé qu'il suffisait que le foncteur % [sic] ont échoués -> ont échoué, accord fautif du participe
\[
T \longrightarrow T_{\mathrm{o}}
\]
soit ou bien lisse et coasphérique, ou bien propre et asphérique - donc tout est OK si $T = T_{\mathrm{o}}$, i.e. si $T$ est ordonnée. Mais dans le cas $T = \Ty$, le foncteur $T \longrightarrow T_{\mathrm{o}}$ n'est ni propre, ni lisse (pas même 0-propre ou 0-lisse).

Finalement j'en suis arrivé à la conviction que dans le cas qui nous occupe ($T = \Ty$, $C = \Ch_{\Sigma}(X;a,b)$), la flèche $C_{\infty} \longrightarrow C$ n'est en général \textit{pas} un isomorphisme dans Hot. En effet, dans le cas $a = b$, si cela \texttt{[était vrai]},
\msp{page 113}
on aurait une description du Hot-groupe $\Omega(X,a) = \Ch(X,a)$ (\footnote[23]{On fait ici $\Sigma = \mathrm{Fl}(X)$.}) en termes d'une catégorie $C$ munie d'une loi de composition \textit{associative} (et unitaire) \textit{dans} Cat, redonnant la loi de groupe de $\Omega$ en passant de (Cat) à Hot. Prenant le nerf de $C$, on trouverait un monoïde semi-simplicial qui, dans Hot, donnerait la loi de groupe de $\Omega(X,a)$. Or je me rappelle (\footnote[24]{C'est un erreur - ça doit être pour d'\textit{autres} objets groupe de Hot que les groupes de lacets. J'ai vu dans Quillen que $\Omega(X,x)$ doit se décrire par un groupe semi-simplicial (quand $X$ est un ensemble semi-simplicial, $x$ un sommet). Donc le problème si $C_{\infty} \simeq C$ dans Hot reste entier.}) qu'on m'a expliqué qu'en général, la loi de groupe de $\Omega(X,a)$ ne pourrait pas être obtenue par une loi \textit{associative} sur un ensemble semi-simplicial. Ainsi, l'insuccès de mes efforts semble tenir, non à ma maladresse, mais à la nature des choses. % [sic] C'est un erreur (note 24) -> C'est une erreur, genre

Je vais résumer les faits positifs auxquels je suis arrivé dans le diagramme suivant, qui a lieu dans la situation générale $C \longrightarrow T$ envisagée au début du présent paragraphe :

\msp{page 114}
% foot: Version 22 décembre 2001
\printedpage{72}
% head: A. Grothendieck : Dérivateurs, Ch. VII, Édition des Universités de Montpellier II et Paris VII
\[
\begin{tikzcd}[column sep=small, row sep=huge]
 & & C_{t'} \arrow[d, "{u^{*}\quad\lab{(pour $u : t \longrightarrow t'$\\ dans $T$)}}"] \arrow[ldd, dotted] \arrow[rdd, dotted] & \\
 & & C_{t} \arrow[ld, swap, "{\lab{flèche définie dans\\ Ind\,Cat\,hot\\ (\texttt{[ou dans]}\\ IndHot,\\ Ind\,($\Hotz$))}}"] \arrow[rd, "{\lab[r]{flèche dans Cat\\ (donc dans\\ Cat\,hot,\\ Hot,\\ $\Hotz$)}}"] & \\
\varinjlim C_{t} = & |[draw, circle]| H_{T} \arrow[rd, swap, "{\lab{flèche\\ définie dans\\ Ind\,Cat\,hot\\ ou dans Ind\,Hot\\ ou dans Ind\,$\Hotz$}}"] & \lab{diagramme\\ commutatif\\ dans Ind\,Hot\\ ou Ind\,($\Hotz$)} & C \\
\lab{dans Ind\,Cat\,hot\\ (ou dans Ind\,Hot\\ ou dans Ind\,($\Hotz$))} & & |[draw, circle]| C_{\infty} \arrow[ru, swap, "{\lab[r]{flèche\\ définie\\ dans Hot,\\ donc dans $\Hotz$,\\ en général pas un\\ isomorphisme, même\\ dans $\Hotz$}}"] & \\
 & & = \varinjlim H_{T} \text{ dans } \Hotz &
\end{tikzcd}
\]

\texttt{[Les annotations dans ce diagramme concernent les flèches en traits pleins.]}

Les objets importants sont $H_{T}$ et $C_{\infty}$, alors que $C$ (sur $T$) et les $C_{t}$ prennent un rôle technique intermédiaire, pour construire $H_{T}$ et $C_{\infty}$. On peut considérer $C_{\infty}$ comme une catégorie (objet de (Cat)), mais dont la construction dépend de nombreux arbitraires, à l'inverse de $H_{T}$ dans Ind\,Cat\,hot. C'est l'objet de $\Hotz$ déduit de $C_{\infty}$ qui a un sens canonique, et peut être décrit comme la limite inductive de $H_{T}$ comme objet de Ind\,$\Hotz$, i.e. comme représentant le même foncteur sur $\mathrm{Hot}_{0}$ que $H_{T}$.

\msp{page 115}
Dans le cas géométrique $T = \Ty$, et $C$ de la forme $\Ch_{\Sigma}(X;a,b)$, donnant lieu à $C_{\infty}(a,b)$, on \textit{n'a pas} de définition naturelle d'un bifoncteur de composition des chemins
\[
C_{\infty}(a,b) \times C_{\infty}(b,c) \ \longrightarrow\ C_{\infty}(a,c).
\]
On peut définir, certes, des accouplements
\[
\Ch_{\Sigma,\tau}(a,b) \times \Ch_{\Sigma,\tau'}(b,c) \ \longrightarrow\ \Ch_{\Sigma,\tau\tau'}(a,c),
\]
ou, avec les notations relatives à la donnée de
\[
\begin{aligned}
&\underline{t} : \mathbf{N}^{o} \longrightarrow \Ty,\\
&\Ch_{n}(a,b) \times \Ch_{m}(b,c) \longrightarrow \Ch_{n+m}(a,c),
\end{aligned}
\]
% foot: Version 22 décembre 2001
\printedpage{73}
% head: A. Grothendieck : Dérivateurs, Ch. VII, Édition des Universités de Montpellier II et Paris VII
mais ces accouplements \textit{ne sont pas} compatibles avec les morphismes de transition des trois systèmes inductifs de catégories $(\Ch_{n})_{n\in\mathbf{N}}$, ils ne le sont \textit{qu'à homotopie près.} Cela suffit cependant pour avoir des accouplements parfaits
\[
H_{T}(a,b) \times H_{T}(b,c) \ \longrightarrow\ H_{T}(a,c),
\]
avec associativité, et (si $\Sigma = \mathrm{Fl}(X)$ tout entier) existence de \textit{vrais} inverses, i.e. de flèches (dans Ind\,Cat\,hot) de `passage au chemin opposé'

\msp{page 116}
\[
H_{T}(a,b) \ \longrightarrow\ H_{T}(b,a), \qquad c \longmapsto c^{o}, \tag*{$(*)$}
\]
donnant lieu au diagramme commutatif exprimant que c'est bien un inverse pour la composition. Une autre façon de le dire, c'est que pour tout objet $Y$ de Cat\,hot (ou de Ind\,Cat\,hot), les ensembles
\[
\mathrm{Hom}_{\mathrm{Cat\,hot}}(Y, H_{T}(a,b)) \overset{\text{déf}}{=} H_{T}^{Y}(a,b)
\]
forment, pour $a,b \in X$ variables, les flèches d'un \textit{groupoïde} (\footnote[25]{Si $\Sigma = \mathrm{Fl}(X)$. Si $\Sigma$ est plus petit, au lieu d'un groupoïde, on trouve une \textit{catégorie} sans plus.}) $\mathrm{Ch}_{\Sigma}^{Y}(X)$ ayant pour ensemble de sommets $\mathrm{Ob}\,X$, et $H_{T}^{Y}(a,b)$ étant formé des flèches de $a$ vers $b$. L'application
\[
H_{T}^{Y}(a,b) \longrightarrow H_{T}^{Y}(b,a), \qquad c \longmapsto c^{o},
\]
déduite du foncteur $(*)$ plus haut n'est autre que l'application de passage à l'inverse. Cette catégorie n'est d'ailleurs autre que celle déduit de % [sic] celle déduit -> celle déduite, accord
\[
X(Y)\Sigma_{Y}^{-1}, \qquad \text{où } \left\{
\begin{aligned}
&X(Y) = \underline{\mathrm{Hom}}_{\mathrm{Cat}}(Y,X)\\
&\Sigma_{Y} \subseteq \mathrm{Fl}(X(Y)),\\
&\text{défini en termes de } \Sigma \text{ `argument par argument'},
\end{aligned}
\right.
\]
à l'aide du foncteur
\[
X \longrightarrow X(Y), \qquad a \longmapsto a_{Y},
\]
\msp{page 117}
associant à tout $a \in X$ le foncteur \textit{constant} sur $Y$ de valeur $a$, en prenant
\[
H_{T}(a,b) \simeq \mathrm{Hom}_{X(Y)\Sigma_{Y}^{-1}}(a_{Y}, b_{Y}) \ .
\]

Dans le même cas géométrique, je me suis convaincu d'ailleurs que $C_{\infty}$ était, en tant qu'objet de Hot (pas seulement de $\Hotz$) \textit{unique à} Hot-\textit{isomorphisme près,} mais (chose étrange) peut-être pas à Hot-isomorphisme \textit{unique} près. De façon générale, revenant aux
% foot: Version 22 décembre 2001
\end{document}
